\documentclass[11pt,reqno]{amsart}
\usepackage[T1]{fontenc}
\usepackage{lmodern,amsmath,amssymb,mathtools,microtype,booktabs,array,longtable}
\usepackage[margin=1in]{geometry}
\usepackage{graphicx}
\usepackage{xr-hyper}
\usepackage[hidelinks]{hyperref}
\ifdefined\pdfinfoomitdate\pdfinfoomitdate=1\fi
\ifdefined\pdftrailerid\pdftrailerid{}\fi
\ifdefined\pdfsuppressptexinfo\pdfsuppressptexinfo=15\fi
\numberwithin{equation}{section}
\newtheorem{theorem}{Theorem}[section]
\newtheorem{proposition}[theorem]{Proposition}
\newtheorem{lemma}[theorem]{Lemma}
\newtheorem{corollary}[theorem]{Corollary}
\theoremstyle{definition}\newtheorem{definition}[theorem]{Definition}
\newtheorem{example}[theorem]{Example}
\theoremstyle{remark}\newtheorem{remark}[theorem]{Remark}
\newcommand{\E}{\mathbb E}\newcommand{\Prb}{\mathbb P}
\newcommand{\R}{\mathbb R}\newcommand{\one}{\mathbf1}
\newcommand{\Q}{\mathcal Q}\newcommand{\TV}{\mathrm{TV}}
\newcommand{\cp}{\mathrm{cp}}\newcommand{\ev}{\mathrm{ev}}\newcommand{\st}{\mathrm{st}}
\DeclareMathOperator{\tr}{tr}\DeclareMathOperator{\dist}{dist}
\DeclareMathOperator{\law}{Law}\DeclareMathOperator{\rank}{rank}
\DeclareMathOperator{\supp}{supp}
\emergencystretch=2em

\externaldocument[][nocite]{supplement-refs}[General_Theta_Foundations_I_R75_Supplement.pdf]
\hypersetup{pdftitle={General Theta Foundations I: Acquired Geometry and Causal Resource Transfer},pdfauthor={Qian Qi}}
\title[Acquired geometry and causal resource transfer]{General Theta Foundations I:\\Acquired Geometry and Causal Resource Transfer}
\author{Qian Qi}\date{October 2026}
\subjclass[2020]{Primary 60J10, 68Q45; Secondary 20C15, 52A20, 60G40, 62C20, 81P15, 90C22, 94A29}
\keywords{causal experiment, predictive state, positive realization, minimal memory, finite group, quantum instrument, convex duality}
\begin{document}
\begin{abstract}
We study the classical memory required to reproduce complete causal
response laws. For a finite observable group, we express the exact
small-contrast memory through irreducible multiplicities in permutation
modules. On a fixed transitive action, the binary response set is
convex if and only if its stabilizer is normal. A permutation-character
witness detects every failure of normality. For every $n\ge5$, we determine
the complete contrast law of a binary $S_n$ response: its minimum
pure-start memory is $n(n-1)+1$ up to and including
$t=(\sqrt n-1)/(\sqrt n+1)$, and $n^2+1$ above that value, although its
observable group and linear dimension remain unchanged. A sharp extremum
for nonnegative singular matrices also classifies the minimum permutation
cores at the transition. For arbitrary
stationary finite-alphabet processes, we prove that a counted pure
initialization costs exactly one state beyond the smallest mixed-start
realization unless the process is independent and identically distributed.
We apply this principle to a qubit instrument with uniformly positive
report probabilities and an exactly depolarizing average channel.
One invariant polygon realizes all
its observed branches and, under an explicit noise bound, all commanded
angles simultaneously. Relative spectral bounds give matching memory
counts in the first sector and in the worst phase. A single finite
simulator also approximates all adaptive stopped laws for policies with
bounded target expected calls, with an explicit resource bound. The complete
causal comparison, optimal discrimination, reset, variational, dual,
spectral and geometric results are retained with their individual
hypotheses and proofs.
\end{abstract}
\maketitle
\section{Introduction}\label{sec:current-overview}\label{sec:introduction}
A predictive state records everything about a past experiment that can
affect its future laws. Its linear dimension does not, by itself, determine
the memory needed to realize those laws with nonnegative probabilities.
Commands can change the next observation, an observation changes the state
on which later commands act, and retaining or communicating information
uses resources. This paper studies the resulting positive-realization
problem within a theory of causal resource comparison and optimal
statistical discrimination.

We distinguish two response objects. A terminal response assigns a
probability row $b(w)$ to each finite command word. An instrument response
assigns a joint probability to each finite command--outcome history. An
exact classical realization uses one fixed collection of stochastic tables
on finitely many live labels. Except in an explicitly stated mixed-start
comparison, the initial label is pure and counted. All lower bounds allow
arbitrary real probabilities. Program-table length, arithmetic precision,
working storage and the protected physical system are separate resources.

The central questions are structural. Which linear predictive actions
admit finite positive realizations? When does linear representation theory
determine the exact resource, and when do probability constraints add a
strict cost? How does the preparation contribute to that cost? Can one
controller reproduce a whole commanded instrument, rather than a scalar
subsequence of its observations? The results below answer these questions
on distinct, explicitly specified classes.

\subsection*{Reversible responses and finite-group structure}
Let $V_b$ be the real span of residual functions $b^w(z)=b(wz)$.
If command shifts are invertible on $V_b$ and generate a group with compact
closure $G_b$, Theorem~\ref{thm:r72-reversible} characterizes finite exact
classical realizability by finiteness of $G_b$. This applies to prepared
finite-dimensional unitary commands after invisible directions have been
removed. Theorem~\ref{thm:r73-normal-form} gives more than a finiteness
criterion: every finite stochastic realization reduces, without increasing
its size, to a permutation core and at most one counted initializer.
The original hidden transition matrices need not form a group. If $N$ is
the number of distinct residuals and $\kappa_G(b)$ is the least size of a
probability factorization $b(g)=\lambda P_gD$, then
\begin{equation}\label{eq:r75-intro-core}
 m(b)=\min\{N,1+\kappa_G(b)\}.
\end{equation}
Here both the preparation $\lambda$ and every decoder row are probability
vectors. Their positivity constraints cannot be replaced by a rank test.

Let $\pi$ be the group average and put $b_t=\pi+t(b-\pi)$.
The retained Theorem~\ref{thm:r74-low-contrast} proves, for every finite
observable group and finite output alphabet, that sufficiently small
positive contrast makes $\kappa_G(b_t)$ equal to the smallest permutation
module containing $V_b$. The new Theorem~\ref{thm:r75-orbit-multiplicity}
turns that minimum into the finite integer problem
\[
 \nu_G(V_b)=\min\left\{\sum_{[H]}q_H[G:H]:
   \sum_{[H]}q_H\dim(\rho^H)\ge r_\rho(b)
   \text{ for every }\rho\in\widehat G\right\}.
\]
The variables are nonnegative integers, $[H]$ ranges over conjugacy classes
of subgroups, and $r_\rho(b)$ is the multiplicity of the complex irreducible
$\rho$ in the complexified response module. Stacked Fourier blocks compute
these multiplicities. The proof includes descent to a real equivariant
embedding, so no real or quaternionic representation type is omitted.
For a prescribed transitive action $H\backslash G$, the same argument gives
the exact number of identical orbit copies needed at small positive
contrast. These are exact resource formulas, without a polynomial-time
claim for subgroup enumeration or integer optimization.

The probability constraint also singles out a precise algebraic boundary.
Theorem~\ref{thm:r75-normality-convexity} proves that the set of binary
responses realizable on one transitive orbit $H\backslash G$ is convex if
and only if $H$ is normal. The forward obstruction is explicit: averaging
the deterministic return responses gives the normalized permutation
character. If $H$ is not normal, one Fourier block has larger rank than
any single-orbit response can have. This obstruction persists at every
positive contrast toward the constant average. When $H$ is normal,
preparation can be absorbed into a quotient decoder, as in the retained
Theorem~\ref{thm:r74-quotient-lp}. Thus the earlier quotient linear programs
have an exact structural explanation. The result concerns this orbit
response set; it does not exclude other linear formulations of special
nonnormal problems.

Probability constraints yield a further, quantitative rigidity. The sharp
singular-matrix bound in Lemma~\ref{lem:r75-singular-coefficient}, combined
with subgroup-index estimates, gives a complete contrast law for an
unbounded nonabelian family. For $G=S_n$, $n\ge5$, let
$\chi_n(g)$ count fixed points of the natural action and define
\[
 \beta_t(g)=\frac{1-t}{n}+\frac{t}{n}\chi_n(g),\qquad
 b_t(g)=(\beta_t(g),1-\beta_t(g)).
\]
Put $\tau_n=(\sqrt n-1)/(\sqrt n+1)$. Theorem~\ref{thm:r75-symmetric-return}
proves, over all finite stochastic controllers,
\[
 m(b_t)=
 \begin{cases}
 1,&t=0,\\
 n(n-1)+1,&0<t\le\tau_n,\\
 n^2+1,&\tau_n<t\le1.
 \end{cases}
\]
Throughout positive contrast, the group is $S_n$, the residual count is
$n!$, and the predictive dimension is $1+(n-1)^2$. The lower bound covers
every finite stochastic controller through \eqref{eq:r75-intro-core}, not
only a chosen collection of natural orbits. The upper tables are explicit;
two command types and binary readout suffice for every $n$. At the
threshold, the coefficient matrix is unique up to coordinate permutation.
Its minimum permutation-core factorizations have fixed preparation rows
and an explicitly described freedom in their orbit weights and decoder
rescaling. This classifies these optimal cores, without asserting that
all hidden stochastic realizations are unique. At full contrast, the
identity value and group mean give the separate point-return rigidity of
Lemma~\ref{lem:r75-return-rigidity}. A diagonal quantum realization supplies
a physical implementation without a quantum-advantage assertion.
The earlier sharp $S_3$ and four-bit examples and the quotient payoff
duals are retained as separate complete model laws.

\subsection*{Initialization, conditional dynamics and one common simulator}
Preparation has a general cost beyond reversible experiments. For any
stationary process over a finite alphabet, a pure hidden initial state has
the same future-response function as a convex combination of its
one-step successors. Unless the process is independent and identically
distributed, that initial state can be eliminated from any pure-start
realization. The resulting exact identity is
\[
 n_{\mathrm{pure}}=n_{\mathrm{mixed}}+1
\]
for every non-i.i.d. stationary process admitting a finite classical model.
Theorem~\ref{thm:r75-stationary-initializer} proves this by an explicit state elimination, including
preservation of every word law. The i.i.d. case has both minima equal to
one; a process with no finite mixed realization has no finite pure one.
A controlled version applies when the predictive preparation is invariant
under a full-support average of unobserved commands. This explains the
initialization charge by stationarity itself, independently of a special
phase or a minimal-rank assumption.

The positive physical mechanism remains the retained
Theorem~\ref{thm:r68-instrument}. For arbitrary known finite instrument
families on $\mathbb C^h$, a uniform lower bound $\eta I/h$ on every
normalized branch yields a common finite physical dictionary. Branchwise
barycentric expansions give one set of nonnegative tables for all commands,
all visible adaptive histories, terminal POVMs, complete stopped laws and
nonnegative history costs. The theorem also provides sparse branch rows,
several-step variants and its stated state-positive noise exponent.
These are inherited results, with their complete proofs in the supplement.

Average mixing alone has a different effect. The retained
Theorem~\ref{thm:r74-depolarizing-obstruction} gives a stationary,
one-dependent qubit process with Hankel rank three, uniformly positive
outcome probabilities and an exactly depolarizing unconditioned channel.
Its exact classical memory is $q$ with a mixed start and $q+1$ with a counted
pure start for a conditional phase of order $q\ge3$, and is infinite for a
phase of infinite order. Theorem~\ref{thm:r75-noisy-memory} develops the corresponding
noisy \emph{complete instrument}. One regular polygon is simultaneously
invariant under all three observed branches. The reset branch still
prepares a pure state at every noise value, so the construction does not
follow by imposing a positive interior margin on all posteriors, even
after a fixed block of calls. The relative spectral lower
bound and this common construction give exact first-sector memory criteria
and the exact worst-phase count
\[
 \max_\theta n_{\mathrm{mixed}}=K(\alpha)
   :=\max\{3,\lceil\pi/\arccos\alpha\rceil\},\qquad 0<\alpha<1.
\]
The scalar sector and polygon geometry were established in the earlier
spectral theory. The additional conclusion here is that its bounds are
attained for the full observed instrument, with one physical dictionary,
and that initialization adds exactly one label by the general theorem.
The zero-retention endpoint $\alpha=0$ has two mixed and three pure
states and is treated separately.

For $0<\alpha<1$ and any integer $K\ge K(\alpha)$, one $K$-vertex
dictionary works for every rotation command in any supplied finite
alphabet. Its tables reproduce all adaptive command--outcome laws. For
$K=K(\alpha)$, the bound is exact when the alphabet includes a maximizing
phase from Theorem~\ref{thm:r75-noisy-memory}. This quantifier differs from choosing a new simulator
for each scalar phase. Theorem~\ref{thm:r75-stopped-approximation} perturbs only the rotating branch and compensates
the third branch to give one finite simulator for the undamped experiment
whose stopped-record error is at most $(1-\alpha)\mathbb E_1[\tau]/8$.
For policies whose undamped target expected calls are at most $B$,
and error tolerance $\delta$, this gives an
explicit $O(\sqrt{B/\delta})$ label upper bound, uniformly over visible
adaptive stopping and command policies. The approximate experiment need
not retain the target expected-duration bound. It does not establish a matching
approximation lower bound or solve a different minimax-risk frontier.

\subsection*{Priority, dependence and the original causal theory}
Residual automata, stochastic realization, Hankel rank, invariant cones and
observable operator models are classical
\cite{RabinScott59,Heller65,CarlylePaz71,Jaeger00,ThonJaeger15,MonrasWinter14}.
The structure of stochastic matrix groups
\cite{Schwarz64,Flor69,Brown64,Schwarz62}, finite-group representation theory
\cite{Serre77}, symmetric positive factorizations
\cite{GPT13,KPT10} and nonnegative-factorization perturbations
\cite{BocciCarliniRapallo11,GillisGlineur12} supply established ingredients.
We credit multiplicity decomposition, Frobenius reciprocity and Fourier
orthogonality as classical facts, with the steps needed here proved below.
The new group assertions are the exact orbit-convexity boundary, its
positive resource interpretation, and the sharp symmetric-group rigidity
law. They extend the already proved reversible normal form and contrast
construction rather than replacing them.

Invariant-polytope simulation of quantum instruments, including adaptive
composition, has explicit antecedents \cite{MonrasWinter14,Okay26}.
Fanizza, Lumbreras and Winter \cite{Fanizza24} treat phase obstructions,
finite quantum processes without finite classical realizations and noisy
classical-memory lower bounds. One-dependent non-HMM processes also
predate this work \cite{HolroydLiggett16}. Neither phenomenon is claimed
as a first. The present assertions concern the general counted-initialization
identity and the displayed qubit family's exact, simultaneous full-instrument
resource laws and stopped-law approximation. Finite diagnostics below are
checks of identities and constructions, not substitutes for the proofs or
for a complete literature-priority assessment.

The main article retains the full actual-history comparison theorem and
its common paid construction. Its witness is evaluated on the competitor's
actual pre-call histories. The paired supplement retains performed-law
compactness, variational and dual attainment, robust tied optimality,
normalized stopping rigidity, finite-width transfer, rank-constrained
optimal discrimination, finite-order quantum measurement, two-stage reset
control, memory and communication accounts, spectral geometry and strategy
implementation. The independently compiled companion retains absolute
calibration and nonintervening directional sensing.

Each layer keeps its own quantifiers: general causal comparison for its
declared experiment class, exact reversible realization under compact
reversibility, stationary initialization for arbitrary stationary processes,
interior-posterior simulation for its stated instrument class, and closed
model laws for their specified data. Local geometric conclusions remain
local. None completes the separate A1, A2 or further Theta construction
objectives. The current theorem and preservation maps distinguish logical
proof dependencies from retained historical material.

\section{Prepared kernels, acquired task laws and observers}\label{sec:data}
All hidden-state, action and report spaces are standard Borel. A fixed unknown index $\theta\in\Theta$ specifies an actual preparation and positive normalized causal kernels. Positivity means nonnegative measures, not an everywhere strictly positive density. Dependence on the entire physical past may be included in the hidden state. A legal policy sees past visible reports and its admitted randomization, but not $\theta$ or a future audit. Iteration of these kernels defines its finite-prefix laws; the consistent countable path law is their kernel extension.

Each call has a cost $c_i\ge c_*>0$. The decision to make the next call is predictable. In the geometric route, the declared productive-call flag is also predictable; in the primitive-instrument route, productive reports are specified outcomes and can be random. Commitment is permitted only after a productive report, and occurs exactly once almost surely in each model. At that time a specified future audit $Z\in H$ is scored, where $H\subset\R^k$ is nonempty, compact and convex. The prediction is committed before that audit is observed. A prospective joint report--audit kernel is defined at every productive call, including rejected calls. The stopping coin is independent of the audit conditional on the entire current visible report and the past. Write $L=\operatorname{diam}(H)^2$; the case $L=0$ is trivial whenever a legal rule exists.

The core cost $C_\tau$ includes initial preparation charged to the core, failed calls and all productive calls through commitment. A separately executed final audit must be added to the raw budget. This convention is fixed before taking a risk infimum. Physical time, stored program constants and temporary computation are not identified with call count.

\subsection*{Three resource interfaces}
A checkpoint procedure may use its whole precommit history, but must output one element of a fixed dictionary of at most $M$ values. This counts its output width only; it is not a claim of $\log_2 M$ total working bits.

An autonomous observer has an alphabet $S$, $|S|\le M$, and one fixed parameter-independent program. At an intercall boundary its only persistent observer memory is its label $s\in S$. Its next command distribution is a fixed function of $s$. A physical port name, elapsed time, subprogram index, random seed or former report is not an extra free input. During a call its current command, current report and fresh coins are available. Whatever persists to another call must be encoded in $S$. A fixed decoder $r:S\to H$ reads only the committed label.

In the \emph{event-stopped} interface a post-report kernel emits $(s',e)\in S\times\{0,1\}$. An event $e=1$ ends the experiment and retains $s'$ for the decoder; no subsequent command is executed. This public commitment event is not a stored bit used in continuing computation. The same label may be used live on one path and terminal on another. In the \emph{state-stopped} interface a fixed map $h_*:S\to\{0,1\}$ determines halt status. The initial label is live, and type-one labels are globally absorbing terminal labels with no executed command rows. Entry into a terminal label is allowed only after a productive report. Thus live and terminal types are disjoint in this interface, and only in this interface by definition.

Write $R_i^\Theta(N,M)$, $i\in\{\cp,\ev,\st\}$, for the infimum of worst-case expected squared loss over programs with robust almost-sure commitment and $\sup_\theta\E_\theta C_\tau\le N$. An empty legal class has value $+\infty$. A program chosen once at random is not a free persistent seed. The ideal measurable class may have fixed real constants; effective implementation is stated separately in Section~\ref{sec:composition}.

\subsection*{A physically acquired comparison law}
Execute a declared prepared comparison experiment and obtain its joint report--audit probability $K_0(dx,dz)$ on $X\times H$. Its report marginal is $\nu$. The law is genuinely preparable and its specification has its own resource account. It need not be a member of $\Theta$. In particular no finite set of calibration data is silently promoted to a uniform certificate. Set
\begin{equation}\label{eq:derived}
 f(x)=\int zK_0(dz\mid x),\quad
 w(x)=\int\|z-f(x)\|^2K_0(dz\mid x),\quad
 d_r(x)=w(x)+\|f(x)-r\|^2.
\end{equation}
Disintegration gives Borel versions with $f\in H$; null sets may be completed arbitrarily. The profile is
\begin{equation}\label{eq:profile}
 \Q_m(b)=\min_{\substack{C\subset H,\ 1\le|C|\le m\\
                  0\le h\le b,\ \int h\,d\nu=1}}
        \int [w+\dist(f,C)^2]h\,d\nu,
 \qquad b\ge1,\ m\ge1.
\end{equation}
We use $+\infty$ for an infeasible cap or width. The quantities $f,w$ are task coordinates of acquired predictive information, not a complete substitute for the predictive quotient. Their derivation and the quotient construction in Appendix~\ref{sec:quotient} precede any use of local dimension.

\subsection*{Lower certificates on actual histories}\label{sec:lower-cert}
Fix \emph{one} actual model $\theta_*\in\Theta$, independent of the competing program. At each actual reached productive pre-call history $h$ and legal command $a$, write $P_{h,a}(dy,dz)$ for its conditional full current report--audit law. Let $P^0_{h,a}$ be a measurable comparison row on the same report space, with a Borel mark $x=\pi_{h,a}(y)$ of law $\nu$ and
\begin{equation}\label{eq:full-moments}
 \E^0[Z\mid y]=f(\pi_{h,a}y),\qquad
 \E^0[\|Z\|^2\mid y]=w(\pi_{h,a}y)+\|f(\pi_{h,a}y)\|^2.
\end{equation}
These are conditional identities given the \emph{full} report, not just its mark. Versions and inequalities are required on actual reachable histories, not only almost everywhere under an unrelated comparison path law.

Fix $0<l_-\le1\le u_-<\infty$. For all Borel $q:Y\to[0,1]$ and all finite-image Borel $r:Y\to H$, suppose predictable nonnegative defects $a_-(h,a),e_-(h,a)$ obey
\begin{align}
 \int q\,d(P_{h,a})_Y
 &\le u_-\int q\,d(P^0_{h,a})_Y+a_-(h,a),\label{eq:lower-mass}\\
 \int q(y)\|z-r(y)\|^2P_{h,a}(dy,dz)
 &\ge l_-\int q(y)\|z-r(y)\|^2P^0_{h,a}(dy,dz)-e_-(h,a).
 \label{eq:lower-task}
\end{align}
Independent current randomization can be adjoined to both rows: integrating inequalities for each fixed random value proves the same inequalities with the same defects. All command-conditioned lower inequalities hold before averaging the chosen command.

Let $A_\tau$ count productive calls. At this same actual model require, for every legal competing procedure,
\begin{equation}\label{eq:effort}
 \E_* C_\tau\ge\kappa_-\E_* A_\tau,\qquad\kappa_->0.
\end{equation}
At budget $N$ suppose, uniformly over all robustly feasible competitors,
\begin{equation}\label{eq:occupation}
 \E_*\sum_{j\le\tau:\ j\ {\rm productive}}a_-(H_{j-1},a_j)\le A_N,
 \qquad
 \E_*\sum_{j\le\tau:\ j\ {\rm productive}}e_-(H_{j-1},a_j)\le E_N.
\end{equation}
These are actual occupation certificates, not properties merely of the proposed upper program. Uniform row defects give $A_N\le(N/\kappa_-)\bar a_-$ and $E_N\le(N/\kappa_-)\bar e_-$. Lemma~\ref{lem:drift} supplies a nonuniform alternative.

\subsection*{One common upper construction}\label{sec:upper-cert}
Independently, let $\mathcal G$ be a parameter-independent program which, from each of its start histories, reaches one productive report. Its conditional expected attempt cost is at most $\kappa_+>0$ on every surviving history in every $\theta\in\Theta$. It uses $B$ live roles and returns to its start label on rejection; no reset of the physical state is assumed. It forms a virtual mark $x$, a current choice $J$ and a permitted export $T_Jr$ of its task readout. For every finite dictionary $C$ assume
\begin{equation}\label{eq:exports}
 |\{T_jr:j\in J,\ r\in C\}|\le D_0|C|,
 \qquad D_0\in\mathbb N.
\end{equation}
Every continuing phase and coin needed for this construction is included in the $B$ live roles. For fixed $0<l_+\le1\le u_+<\infty$ and $a_+,e_+\ge0$, assume for all Borel $q:X\to[0,1]$ and finite-image Borel $r:X\to H$, conditionally on every such history,
\begin{align}
 \E[q(x)]&\ge l_+\int q\,d\nu-a_+,\label{eq:upper-mass}\\
 \E[q(x)\|Z-T_Jr(x)\|^2]
 &\le u_+\int q(x)d_{r(x)}(x)\nu(dx)+e_+.
 \label{eq:upper-task}
\end{align}
The expectation here includes the within-attempt choice $J$. Unlike the lower, the upper is not required to hold for each choice separately. The choices retained after this call are nevertheless charged through the actual export or live role. These hypotheses express one common causal construction, not modelwise optimizers.

\section{Reversible response groups and common controlled memory}\label{sec:r72-reversible}
The fixed-phase results concern one transition matrix. Several commands require one family of matrices on the same retained labels, and checking every command in isolation is insufficient. We give an intrinsic finite-realization criterion for responses whose minimal linear dynamics is reversible and compact, including any finite family of unitary commands on a prepared finite-dimensional quantum system. We then obtain a matching arithmetic memory law for a pair of noncommuting involutions. This is a theorem about exact controlled response realization, not an assertion that a simulator's private label is an observed sufficient statistic for arbitrary feedback optimization.

Positive realization through invariant polyhedral sets is classical \cite{Heller65,MonrasWinter14,ThonJaeger15}. The argument below uses a section of the realizing simplex before passing to the observable quotient; auxiliary hidden-state responses need not belong to the minimal response space. Compact reversibility then turns containment into equality. The resulting finite-group obstruction concerns all command words, rather than a single spectral restriction.

\subsection{The minimal response space}
Let $A$ and $C$ be nonempty finite alphabets of commands and terminal outcomes. A supplied response is a function $b:A^*\to\Delta(C)$. Its residual after a word $w$ is $b^w(z)=b(wz)$, including the empty word $\varnothing$. Put
\[
 V_b=\operatorname{span}_{\R}\{b^w:w\in A^*\},
 \qquad (L_af)(z)=f(az).
\]
Assume $V_b$ is finite dimensional. It is invariant under each $L_a$, since $L_ab^w=b^{wa}$. Call the response \emph{compactly reversible} if the restrictions $L_a|_{V_b}$ are invertible and the closure $G_b$ of their generated group is a compact subgroup of $\operatorname{GL}(V_b)$. This concerns the observable quotient; it does not require the full physical command group to be finite or fully observable.

An exact stochastic realization with $M$ live labels has row-stochastic matrices $Q_a$, one row-stochastic terminal decoder matrix $D$, and the prescribed initial row $e_1^{\mathsf T}$, such that
\begin{equation}\label{eq:r72-response}
 b(a_1\cdots a_n)=e_1^{\mathsf T}Q_{a_1}\cdots Q_{a_n}D.
\end{equation}
Randomized terminal readout is allowed even when all $Q_a$ are deterministic. The target input word is not encoded into an uncounted clock or persistent seed.

\begin{theorem}[Finite realization of compactly reversible responses]\label{thm:r72-reversible}
For a compactly reversible response, the following are equivalent:

\textup{(i)} An exact finite stochastic realization exists.

\textup{(ii)} The residual set $\{b^w:w\in A^*\}$ is finite.

\textup{(iii)} The observable group $G_b$ is finite.

\textup{(iv)} An exact realization with deterministic transitions and a possibly randomized decoder exists.

The least number of deterministic-transition labels is the residual cardinality. This is not in general the least stochastic label count. If an $M$-label stochastic realization exists, then
\[
 |G_b|\le(2^M-1)!.
\]
This last bound is only a finite combinatorial certificate, not a sharp memory law.
\end{theorem}
\begin{proof}
Suppose \eqref{eq:r72-response} holds, and write $r_w=e_1^{\mathsf T}Q_w$. Let $R=\operatorname{span}_{\R}\{r_w\}\subseteq\R^M$. The assignment $\pi(r_w)=b^w$ extends to a well-defined surjective linear map $\pi:R\to V_b$: a relation $\sum_wc_wr_w=0$ implies, after right multiplication by $Q_zD$, the identical relation among $b^w(z)$ for every $z$. Moreover
\[
 \pi(rQ_a)=L_a\pi(r)\qquad(r\in R).
\]
The simplex section $K=R\cap\Delta_M$ is a compact polytope, invariant under right multiplication by every $Q_a$. Its image $P=\pi(K)$ is a compact polytope, contains every residual, and spans $V_b$. It satisfies $L_aP\subseteq P$. It is important that we intersect the simplex with $R$ first: projecting arbitrary hidden-state responses directly into $V_b$ would not be justified.

For an element $T$ of a compact matrix group, positive powers have a subsequence tending to the identity with exponents tending to infinity. One elementary proof takes a convergent subsequence of powers with successively growing index gaps; quotients of two close powers approach the identity. A finite-order element is covered by its period multiples. Thus positive powers also approach $T^{-1}$. Applying this to $L_a$ and using closedness of $P$ gives $L_a^{-1}P\subseteq P$, and hence $L_aP=P$. All elements of $G_b$ preserve $P$. They permute its finitely many vertices, and this permutation action is faithful because those vertices span $V_b$. Consequently $G_b$ is finite. This proves (i)$\Rightarrow$(iii).

For the displayed bound, a vertex of $K$ is determined by its positive support. If two distinct vertices had the same support, their difference could be added and subtracted in sufficiently small multiples while preserving $R$, unit mass and nonnegativity, contradicting extremality. Thus $K$ has at most $2^M-1$ vertices. A linear image of a polytope has no more vertices than the original polytope, so the faithful vertex action proves the bound.

If $G_b$ is finite, the residual orbit is finite. Conversely, if the residual set is finite, each invertible $L_a$ maps it injectively into itself, hence permutes it. Since it spans $V_b$, the generated group and its closure are finite. This proves (ii)$\Leftrightarrow$(iii). Under (ii), use the residuals as labels, send $b^w$ to $b^{wa}$ on input $a$, and use decoder $b^w(\varnothing)$. This realizes $b$ with deterministic transitions. If any such deterministic-transition machine sends $u$ and $v$ to the same label, every suffix yields the same decoder law, so $b^u=b^v$. This proves its minimal count. Finally (iv)$\Rightarrow$(i) is immediate. No assertion that stochastic transitions can always be determinized without increasing size was used.
\end{proof}

\begin{corollary}[Prepared unitary commands]\label{cor:r72-unitary}
Let $U_a$ be finitely many unitary matrices on $\mathbb C^d$, let $\rho$ be a prepared density matrix, and let $\{E_r:r\in C\}$ be a terminal positive-operator-valued measurement. For
\[
 b(w,r)=\tr(E_r\Phi_w(\rho)),\qquad \Phi_a(X)=U_aXU_a^*,
\]
exact finite stochastic realization for all command words is possible if and only if its observable response group is finite. The group is computed after quotienting only those directions invisible to every future commanded terminal test. Its dimension is at most $d^2$.
\end{corollary}
\begin{proof}
Let $W=\operatorname{span}_{\R}\{\Phi_w(\rho)\}$ in the Hermitian matrices. Each conjugation is a Hilbert--Schmidt isometry. Its forward invariant subspace $W$ is also backward invariant, by finite-dimensional invertibility. Let
\[
 N=\{X\in W:\tr(E_r\Phi_zX)=0\text{ for all }z,r\}.
\]
It is invariant under every $\Phi_a$ and hence under its inverse. The orthogonal complement of $N$ in $W$ is invariant as well. Therefore the induced transformations on $W/N$ are orthogonal for the quotient inner product. The map $X\mapsto(z,r\mapsto\tr(E_r\Phi_zX))$ identifies $W/N$ with $V_b$ and intertwines them with $L_a$. Their generated group has compact closure in the orthogonal group. Apply Theorem~\ref{thm:r72-reversible}. A constant measurement can make this induced group trivial even when the physical unitary group is infinite, so finiteness of the full physical group is not a necessary condition.
\end{proof}

\subsection{A matching memory law for two noncommuting commands}
Fix $x\in\R$ with $\sin(2\pi x)\ne0$. On a protected qubit take
\[
 U_0=\begin{pmatrix}1&0\\0&-1\end{pmatrix},\qquad
 U_1=\begin{pmatrix}\cos(\pi x)&\sin(\pi x)\\
                     \sin(\pi x)&-\cos(\pi x)\end{pmatrix},
 \quad \rho=E=|0\rangle\langle0|.
\]
Both commands are involutions. For known $0\le\alpha\le1$, their channels are
\[
 \Phi_{a,\alpha}(X)=\alpha U_aXU_a^*+(1-\alpha)\tr(X)I/2.
\]
Let $M_x^{\rm ctrl}(\alpha)$ be the least stochastic live memory of the whole binary response $w\mapsto\tr(E\Phi_{w,\alpha}\rho)$, using the same matrices and decoder for every word. In the real Bloch plane, with coordinates $(z,x)$, these commands are $\alpha S_0$ and $\alpha S_1$, where $S_0$ reflects the angle $t$ to $-t$ and $S_1$ reflects it to $2\pi x-t$. Their product is $\alpha^2R_{2\pi x}$.

\begin{theorem}[All-word arithmetic law for an involutive command pair]\label{thm:r72-control-envelope}
For $0<\epsilon<1$ and $\alpha=1-\epsilon$, the arithmetic envelope of Section~\ref{sec:r71-arithmetic} gives
\begin{equation}\label{eq:r72-control-envelope}
 \tfrac14\mathfrak A_x(\epsilon)
 \le M_x^{\rm ctrl}(1-\epsilon)
 \le(1+4\pi^2)\mathfrak A_x(\epsilon).
\end{equation}
The upper is one explicit finite stochastic machine valid for every mixed command word, not separately chosen unary machines. The lower ranges over arbitrary common stochastic realizations.

For irrational $x$, the lower and upper logarithmic growth exponents as $\epsilon\downarrow0$ are respectively
\[
 \frac1{\mu(x)+1},\qquad\frac{\mu(x)-1}{2\mu(x)-1},
\]
with values $0,1/2$ at $\mu=\infty$. A uniform cubic-root law holds precisely for badly approximable $x$; otherwise the normalized count has lower limit zero and upper limit infinity.

At $\alpha=1$, the exact count is the reduced denominator $q$ if $x$ is rational, and is infinite otherwise. At $\alpha=0$, it is two. Every strictly damped pair has a finite exact common realization. These statements concern controlled response memory, not a general multi-command boundary-optimal policy theorem.
\end{theorem}
\begin{proof}
Any common realization restricted to the words $(01)^n$ realizes
\[
 \tfrac12[1+\alpha^{2n}\cos(2\pi nx)]
\]
with transition matrix $Q_0Q_1$. Its size is therefore at least
$m_x(1-\alpha^2)\ge\mathfrak A_x(\epsilon(2-\epsilon))/2$
by Theorem~\ref{thm:r71-envelope}. For $1\le c\le2$, the defining maximum of the envelope obeys $\mathfrak A_x(c\epsilon)\ge\mathfrak A_x(\epsilon)/c$ whenever both arguments are in its domain. This follows for each denominator before taking its infimum. Taking $c=2-\epsilon$ proves the lower in \eqref{eq:r72-control-envelope}.

For the upper, the rational multiple-polygon construction in the proof of Theorem~\ref{thm:r71-envelope} gives a regular $l$-gon $P$ on the unit circle, with vertex $v_0=(1,0)$, such that
\[
 l\le(1+4\pi^2)\mathfrak A_x(\epsilon),\qquad
 \alpha R_{2\pi x}v_j\in P\quad\text{for every vertex }v_j.
\]
The first reflection $S_0$ permutes this regular polygon and zero belongs to it. Consequently $\alpha S_0P\subseteq P$, while $S_1=R_{2\pi x}S_0$ gives $\alpha S_1P\subseteq P$. Choose barycentric rows $Q_a(j,k)$ of $\alpha S_av_j$ in the same vertices. Let $d_j=(1+(v_j)_1)/2$. Starting at vertex zero, induction on the actual command word shows that the expected retained vertex equals the protected Bloch vector. Pairing with the first coordinate proves every response identity. The start is one of the counted vertices; there is no extra initialization mixture.

The two constants in \eqref{eq:r72-control-envelope} transfer both exponents and both uniform-law statements from Theorems~\ref{thm:r71-exponents} and \ref{thm:r71-cubic-rigidity}. At zero damping, a rational phase of reduced denominator $q\ge3$ has a common deterministic $q$-gon realization. The word-pair transition has a primitive $q$th-root mode and needs at least $q$ states by the unit-circle spectral obstruction. For irrational $x$ that same obstruction rules out every finite size. At $\alpha=0$ the response is one at the empty word and $1/2$ after every nonempty word, realized by a first label followed by an absorbing fair label; one label cannot change its decoder. This proves all endpoints.
\end{proof}

A fully algebraic example uses $U_1=5^{-1/2}\left(\begin{smallmatrix}2&1\\1&-2\end{smallmatrix}\right)$. Then $e^{2\pi ix}=(3+4i)/5$ is not a root of unity, since its sum with its conjugate is the rational noninteger $6/5$. Each undamped command in isolation has period at most two; their common all-word response has no finite stochastic realization. The arithmetic exponent of this particular angle is not determined here. The example isolates the failure of commandwise testing without inferring a Diophantine class from finitely many coefficients.

If an independent herald with fixed probability $p$ is placed before each command operation, and only failures apply the supplied command, every common realization above preserves the terminal response under every external command schedule and the first-herald cost $1/p$. The commands are actual inputs to that same retained-state transducer. This does not make its private label observable to the caller, allow uncharged midstream quantum measurements, or prove optimality for a controller free to choose an easier command sequence. The unary optimal squared-task interpretation below continues to use its separately proved physical embedding.

\section{A permutation normal form with an exact memory charge}
\label{sec:r73-normal-form}

Theorem~\ref{thm:r72-reversible} identifies when finite memory exists. We now
determine the form of every possible finite realization, retaining the cost of
its prescribed pure initial label. The distinction between a minimal linear
representation and a positive realization remains essential. In particular,
the permutation representation carried by the observable space does not by
itself provide a decoder with values in the probability simplex.

Let $b$ be a compactly reversible response admitting a finite realization.
Write $G$ for its finite observable group, $d=\dim V_b$, and $N$ for its residual
cardinality. We use the chronological convention $g(u)g(v)=g(uv)$: the right
action on $V_b$ is $f\cdot g(w)=L_wf$, where $L_wf(z)=f(wz)$. Thus $G$ in this
notation is the opposite of the group of operators with the usual convention
for composition; inversion identifies the two groups. No group cardinality,
faithfulness, or representation-theoretic statement depends on this convention.
Every element has a representative in $A^*$, including its inverse, since a
finite group generated by the command elements is already generated by their
positive words. The value $b(g)=b(w)$ is consequently well defined.

A \emph{probability factorization on a finite $G$-set} consists of a finite set
$X$, row permutation matrices $P_g$ with $P_gP_h=P_{gh}$, a probability row
$\lambda$ on $X$, and a row-stochastic matrix $D$ with columns indexed by $C$,
such that
\begin{equation}\label{eq:r73-probability-factorization}
 b(g)=\lambda P_gD\qquad(g\in G).
\end{equation}
All rows of $D$ are probability laws, including rows associated with states
that are not distinguished by the immediate terminal measurement. Unused
orbits may be deleted. Define
\[
 \kappa_G(b)=\min\{|X|:\text{\eqref{eq:r73-probability-factorization} holds}\}.
\]
The residual realization shows that this minimum is finite and at most $N$.
It permits a mixed $\lambda$ without charging for its preparation; the
stochastic memory $m(b)$ continues to charge every live label, beginning at
the specified point mass as in \eqref{eq:r72-response}.

\begin{theorem}[Reversible stochastic normal form]\label{thm:r73-normal-form}
Every $M$-label stochastic realization of $b$ gives a faithful finite $G$-set
$X$ and a factorization \eqref{eq:r73-probability-factorization} satisfying
\begin{equation}\label{eq:r73-initialization-charge}
 |X|+\mathbf 1_{\{\lambda\text{ is not a point mass}\}}\le M.
\end{equation}
Conversely, every such factorization has a realization with the displayed
number of labels. Its transitions on $X$ are permutations; if $\lambda$ is
mixed, its only additional label initializes the distribution on the first
command. Thus arbitrary continuing stochastic dynamics confers no advantage
over one initialization mixture followed by a permutation action.
\end{theorem}

\begin{proof}
Fix a realization $(e_1^{\mathsf T},\{Q_a\},D_0)$, and retain the reachable
row space $R$ and the surjection $\pi:R\to V_b$ from the proof of
Theorem~\ref{thm:r72-reversible}. Write
\[
 \mathcal L=\{w\in A^*:g(w)=1\}.
\]
These are identity words in the observable group; their matrices need not be
identity matrices on hidden states. Choose finitely many words
$\ell_1,\ldots,\ell_s\in\mathcal L$ whose positive matrix entries cover all
directed edges occurring in any $Q_\ell$, $\ell\in\mathcal L$. There are at
most $M^2$ distinct edges, so $s\le M^2$ is possible. If no nontrivial choice
is needed, the empty word suffices. Set
\[
 T=\frac{I+Q_{\ell_1}+\cdots+Q_{\ell_s}}{s+1}.
\]
Its closed communicating classes $C_1,\ldots,C_k$ are closed for every
$Q_\ell$ with $\ell\in\mathcal L$. The positive diagonal makes the
restriction of $T$ to each such class aperiodic. The elementary convergence
argument in Lemma~\ref{lem:r73-lazy-convergence} gives a stochastic idempotent
\[
 E=\lim_{n\to\infty}T^n,
 \qquad
 \Delta_ME=\operatorname{conv}\{\mu_1,\ldots,\mu_k\},
\]
where $\mu_i$ is the unique stationary probability row of $T$ on $C_i$.
Their supports are disjoint, so this image is a simplex with exactly $k$
vertices. Every identity-word matrix preserves each $C_i$; hence
\begin{equation}\label{eq:r73-sandwich}
 EQ_\ell E=E\qquad(\ell\in\mathcal L).
\end{equation}
This is a sandwiched equality. We do not assert $Q_\ell E=E$, $EQ_\ell=E$,
or common stationarity of the $\mu_i$ under all identity words.

Each $T^n$ is a convex combination of identity-word matrices. Consequently
$RE\subseteq R$ and $\pi(rE)=\pi(r)$ for $r\in R$. For a word $u$, put
$F_u=EQ_uE$, considered as a stochastic map of $\Delta_ME$ to itself.
Choose a word $\bar u$ representing $g(u)^{-1}$. Replacing the middle $E$
by the convex combinations $T^n$ and using
$u\ell\bar u\in\mathcal L$ gives
\[
 F_uF_{\bar u}=E=F_{\bar u}F_u.
\]
Thus $F_u$ is an affine automorphism of that simplex. Such an automorphism
permutes its vertices, because both it and its inverse preserve the simplex
and therefore preserve extremality. If $g(u)=g(v)$, the same inverse word
shows $F_uF_{\bar u}=F_vF_{\bar u}=E$, so $F_u=F_v$. Next,
$EQ_uT^nQ_vE$ is a convex combination of $F_{u\ell v}$, all equal to
$F_{uv}$. Passing to the limit proves $F_uF_v=F_{uv}$. We have obtained a
permutation action $P_g$ on the vertices $X=\{\mu_1,\ldots,\mu_k\}$.

This action is faithful. If $F_u$ is the identity on $\Delta_ME$, then
$rEQ_uE=rE$ for all rows $r$. For $r\in R$, application of $\pi$ shows
$L_u\pi(r)=\pi(r)$. Surjectivity gives $L_u=I$ on $V_b$, and so $g(u)=1$.
Write $e_1^{\mathsf T}E=\sum_i\lambda_i\mu_i$ and set
$D(i,\cdot)=\mu_iD_0$. Intertwining on $R$ gives, for every word $w$,
\[
 \lambda P_{g(w)}D=e_1^{\mathsf T}EQ_wED_0=b(w),
\]
including the empty word. If the original initial label lies in $C_i$, its
row under $E$ is exactly $\mu_i$. A mixed $\lambda$ therefore requires that
the initial label lie outside the union of the closed classes. The $k$
nonempty disjoint classes then use at most $M-1$ labels. This proves
\eqref{eq:r73-initialization-charge}; a pure $\lambda$ only requires $k\le M$.

For the converse, when $\lambda$ is a point mass its support point is the
initial label in $X$. If it is mixed, adjoin a label $\star$ with decoder
$D(\star,\cdot)=\lambda D$. On command $a$, its transition row into $X$ is
$\lambda P_{g(a)}$, and no transition returns to $\star$. The rows on $X$
are those of $P_{g(a)}$. Induction gives \eqref{eq:r73-probability-factorization}
for every nonempty word, and the decoder at $\star$ gives the empty word.
The first input is applied during initialization, so it is not lost or
encoded in an uncounted preparatory clock.
\end{proof}

The general structure of nonnegative idempotents and bounded groups of
nonnegative matrices is classical; see Brown~\cite{Brown64},
Schwarz~\cite{Schwarz64}, and Flor~\cite{Flor69}. The proof above supplies the
additional steps required for an arbitrary semigroup realization: the
identity-word projection is constructed before passing to the observable
quotient, the sandwiched action is proved to factor through that quotient,
and the prescribed pure start is charged. It does not assume that the
original stochastic matrices themselves form a group.

\begin{theorem}[Exact memory formula]\label{thm:r73-memory-formula}
With randomized terminal readout allowed and the initial label prescribed
pure, the minimum stochastic memory is
\begin{equation}\label{eq:r73-memory-formula}
 m(b)=\min_{(X,\lambda,D)}
 \left(|X|+\mathbf1_{\{\lambda\text{ is not a point mass}\}}\right)
 =\min\{N,1+\kappa_G(b)\}.
\end{equation}
The minimum ranges over precisely the probability factorizations
\eqref{eq:r73-probability-factorization}, not arbitrary linear factorizations.
\end{theorem}
\begin{proof}
The first equality is both directions of Theorem~\ref{thm:r73-normal-form}.
A point-mass factorization gives a deterministic-transition realization on
$X$, so $|X|\ge N$ by Theorem~\ref{thm:r72-reversible}. A mixed factorization
costs at least $1+\kappa_G(b)$. Conversely, the residual machine has $N$
labels, and a factorization attaining $\kappa_G(b)$ costs at most
$1+\kappa_G(b)$. These inequalities give the second equality, also when a
minimum factorization happens to have a pure start.
\end{proof}

\begin{corollary}[Sharp group bounds]\label{cor:r73-group-bound}
An $M$-label realization of a compactly reversible response satisfies
\[
 |G|\le M!,\qquad
 M<N\ \Longrightarrow\ |G|\le(M-1)!.
\]
Both bounds are attained by families below. In particular, a stochastic
memory saving over the residual count requires a faithful permutation core
on at most $M-1$ points.
\end{corollary}
\begin{proof}
The core action in Theorem~\ref{thm:r73-normal-form} embeds $G$ into $S_k$.
If $M<N$, its start cannot be pure, since that would produce a deterministic
machine with at most $M$ labels. Thus $k\le M-1$ in this case.
\end{proof}

Let $\nu_G(V_b)$ be the least cardinality of a finite $G$-set whose real
permutation module contains $V_b$, and let $\mu(G)$ be the least cardinality
of a faithful finite $G$-set. All references to $V_b$ here use the right-action
convention already fixed. A probability factorization gives $V_b$ as a
subquotient of $\mathbb R^X$: take the cyclic span of $\lambda$, then its
observable quotient. Averaging an inner product over finite $G$ splits every
invariant subspace orthogonally, so this subquotient embeds in
$\mathbb R^X$. We obtain
\begin{equation}\label{eq:r73-module-bounds}
 d\le\nu_G(V_b)\le\kappa_G(b),\qquad
 m(b)\ge\max\{d,\mu(G),\min(N,1+\nu_G(V_b))\}.
\end{equation}
These are lower bounds, not an identification of a linear module with a
bounded decoder. Proposition~\ref{prop:r73-decoder-obstruction} makes the
distinction strict.

\begin{corollary}[Rank rigidity]\label{cor:r73-rank-rigidity}
For a finite compactly reversible response,
\[
 m(b)=d\quad\Longleftrightarrow\quad N=d.
\]
In that case every realization with $d$ labels has permutation transition
matrices. If $N>d$, then $m(b)\ge d+1$.
\end{corollary}
\begin{proof}
Every $M$-label realization has $d\le M$. If $M=d$, the core in
Theorem~\ref{thm:r73-normal-form} has $k\ge d$ by its linear rank, hence
$k=d$ and its start is pure. Thus $N\le d$, whereas $N\ge d$ because the
residuals span $V_b$. Conversely $N=d$ gives a deterministic machine with
$d$ labels. For the assertion about the original matrices, a $d$-label
realization has $R=\mathbb R^d$ and the observable map $\pi$ is an
isomorphism. An inverse word for $g(a)$ therefore gives a stochastic inverse
of $Q_a$ on all of $\mathbb R^d$. A stochastic map with a stochastic inverse
is an affine automorphism of the full simplex and hence a permutation.
\end{proof}

\subsection{Two commands and a factorial deterministic gap}
Let $k\ge3$, let $P_g$ be the natural row permutation action of $S_k$, choose
a nonconstant $\lambda\in\Delta_k$, and choose a nonconstant column
$t\in[0,1]^k$. Two commands are the $k$-cycle $(1\,2\,\cdots\,k)$ and the
transposition $(1\,2)$; these generate $S_k$. Supply the binary response
\begin{equation}\label{eq:r73-symmetric-response}
 b(g)=\bigl(\lambda P_gt,\,1-\lambda P_gt\bigr).
\end{equation}
If the distinct coordinate values of $\lambda$ have multiplicities
$n_1,\ldots,n_j$, put $N_\lambda=k!/(n_1!\cdots n_j!)$.

\begin{proposition}[Exact symmetric-group family]\label{prop:r73-symmetric-family}
The response \eqref{eq:r73-symmetric-response} has observable group $S_k$,
dimension $k$, residual cardinality $N_\lambda$, and
\[
 \kappa_{S_k}(b)=k,\qquad
 m(b)=\min\{N_\lambda,k+1\}.
\]
When the coordinates of $\lambda$ are pairwise distinct, the least stochastic
memory is $k+1$ and the least deterministic-transition memory is $k!$.
For a point-mass $\lambda$, both are $k$. If either $\lambda$ or $t$ is
constant, the response instead has one residual and one-label memory.
\end{proposition}
\begin{proof}
For any nonconstant vector, differences between it and its transposes yield
a nonzero multiple of some $e_i-e_j$. Conjugating by all permutations shows
that the span of its orbit differences is
$\mathbf1^\perp$. The orbit of $\lambda$ therefore spans $\mathbb R^k$,
because every row has mass one. The orbit of $t$ also spans
$\mathbb R^k$: its differences span $\mathbf1^\perp$, and its coordinate
sum is positive, since a nonconstant nonnegative vector is nonzero.
Consequently the map
\[
 r\longmapsto\bigl(h\longmapsto(rP_ht,\,r\mathbf1-rP_ht)\bigr)
\]
is injective on $\mathbb R^k$. It identifies the reachable row space with
$V_b$, and identifies the observable action with the faithful natural
permutation action. Its residuals are precisely the distinct permutations
of $\lambda$, giving $d=k$ and $N=N_\lambda$. The displayed factorization
and the rank lower bound give $\kappa=k$. The memory formula proves the
remaining assertion; it also shows $N_\lambda\ge k$. The two constant
cases are immediate from \eqref{eq:r73-symmetric-response}.
\end{proof}

This is an exact prepared quantum experiment in every dimension $k$: take
$\rho=\operatorname{diag}(\lambda)$, permutation unitaries for the commands,
and terminal effect $E=\operatorname{diag}(t)$. The construction already
lies in the diagonal, classical subalgebra. It demonstrates a stochastic
versus deterministic memory separation, and does not assert a quantum
advantage for this family. It also proves sharpness of
Corollary~\ref{cor:r73-group-bound}: the point-mass family has $M=k$ and
$|G|=k!$, while pairwise distinct coordinates have $M=k+1<N=k!$ for $k\ge3$.

For a concrete rational example, take $k=3$,
$\lambda=(1/2,1/3,1/6)$, and $t=(0,1/4,1)^{\mathsf T}$.
The six first-outcome values are
\[
 \left\{\frac14,\frac38,\frac7{24},\frac{13}{24},
                 \frac{11}{24},\frac7{12}\right\}.
\]
It has $d=3$, $N=6$, $\kappa=3$, and $m=4$.

\subsection{The decoder constraint can remove the entire saving}
The representation-theoretic lower bound in
\eqref{eq:r73-module-bounds} need not attain the positive factorization
minimum, even for the same group, rank, and residual count as this example.

\begin{proposition}[A strict bounded-decoder obstruction]
\label{prop:r73-decoder-obstruction}
Let $u=(1,-1,0)$ and let $S_3$ act by its natural permutation matrices. The
binary response with first coordinate
\begin{equation}\label{eq:r73-endpoint-response}
 \beta(g)=\frac12+\frac14 uP_gu^{\mathsf T}
\end{equation}
has $d=3$, $N=6$, and $\nu_{S_3}(V_b)=3$, but
\[
 \kappa_{S_3}(b)=m(b)=6.
\]
It is realized by a prepared qubit with
\[
 U_0=\begin{pmatrix}1&0\\0&-1\end{pmatrix},\qquad
 U_1=\frac12\begin{pmatrix}1&\sqrt3\\\sqrt3&-1\end{pmatrix},
 \qquad \rho=E=|+\rangle\langle+|.
\]
Thus the least stochastic memory can differ for responses having the same
observable group, module, and residual cardinality.
\end{proposition}
\begin{proof}
The orbit of $u$ consists of its six distinct permutations and spans the
two-dimensional zero-sum subspace. Pairing with all permutations of $u$
is nondegenerate on that subspace. Including the constant outcome sum gives
dimension three and six distinct residuals. This representation is the
natural permutation module $\mathbb R^3$, and the rank bound gives
$\nu=3$.

Write $s=(1\,2)$ and $r=(1\,2\,3)$. Directly from
\eqref{eq:r73-endpoint-response},
\[
 \beta(1)=1,\quad\beta(s)=0,\quad
 \beta(r)=\beta(r^2)=\frac14,
 \quad\beta(\tau)=\frac34
 \quad\text{for the other two transpositions }\tau.
\]
Suppose a probability factorization exists on at most five points. Its
action is faithful, by the same residual argument as in
Theorem~\ref{thm:r73-normal-form}. Transitive $S_3$-sets have sizes
$1,2,3,6$, since their stabilizers have orders $6,3,2,1$; the subgroups of
order two are conjugate and the subgroup of order three is unique. Thus
our faithful action has one natural three-point orbit and, at most, one
two-point sign orbit together with fixed points. Two three-point orbits
already exceed the budget, and actions on only fixed and sign orbits
annihilate the alternating subgroup.

Let $a$ be any point of positive initial mass and write $q$ for the
first-coordinate decoder. The endpoint values force $q(a)=1$ and
$q(a\cdot s)=0$. No fixed point can carry mass. Within the natural orbit,
the point fixed by $s$ cannot carry mass, and both of the points exchanged
by $s$ cannot carry positive mass simultaneously. Thus the initial mass in
that orbit is concentrated at one moved point $a$; denote its total by
$A$. If a sign orbit is present its mass is concentrated at a single point;
write this total as $B$. We have $A+B=1$. Every even permutation has decoder
one on the mass in the sign orbit, and every odd permutation has decoder
zero there.

The images $a\cdot r$ and $a\cdot r^2$ are the other two points in the
natural orbit. One is $a\cdot s$, with decoder zero, and the other's decoder
is some $t\in[0,1]$. The two equal values $\beta(r)=\beta(r^2)=1/4$ give
$B=1/4$ and $At=0$. Hence $A=3/4$ and $t=0$. Of the two transpositions
different from $s$, one fixes $a$ and the other sends it to the third point.
Their first-outcome values would therefore be $A$ and zero, whereas both
must be $3/4$. This contradiction rules out all factorizations on at most
five points. The regular six-point residual machine gives $\kappa\le6$,
so $\kappa=6$ and Theorem~\ref{thm:r73-memory-formula} gives $m=6$.

For the physical realization, the adjoint actions of $U_0$ and $U_1$ on
the real Bloch plane are two reflections whose product is a rotation by
$2\pi/3$. The initial and terminal Bloch vector is the $x$-axis vector.
The first reflection sends it to its negative, the other two reflections
have inner product $1/2$ with it, and the two nonidentity rotations have
inner product $-1/2$. The Born probability is one half of one plus that
inner product. It has exactly the displayed six group values, so it is
\eqref{eq:r73-endpoint-response} after the specified generator assignment.
\end{proof}

The cyclic partition formula of the preceding revision is recovered as a
special case in Proposition~\ref{prop:r73-cyclic-module}. For the general
finite group, however, both the start charge and the bounded-decoder
constraint belong in the exact minimization problem. Neither can be
removed by a rank calculation.

\section{Exact memory under terminal readout contrast}
\label{sec:r74-contrast}

The normal form separates the cost of a permutation module from the positivity
of its preparation and decoder. We now determine when that linear lower bound
is attained, for every finite observable group, and describe the exact memory
transitions as the terminal contrast varies. This is a change of terminal
readout. It differs from the length-dependent background damping of
Section~\ref{sec:r73-finite-certificates}.

Let $b$ be a nonconstant compactly reversible response with finite observable
group $G$, residual cardinality $N$, and real observable module $V_b$. Set
\begin{equation}\label{eq:r74-contrast}
 \pi=\frac1{|G|}\sum_{g\in G}b(g),\qquad
 b_t(g)=\pi+t\{b(g)-\pi\},\qquad 0\le t\le1.
\end{equation}
If $\pi_c=0$, nonnegativity implies $b(g,c)=0$ for every $g$; such an output
can be removed and later restored as a zero decoder column. Thus positivity
of all remaining $\pi_c$ imposes no additional hypothesis. For prepared
unitary commands, \eqref{eq:r74-contrast} is realized by the terminal
measurement $E_c(t)=tE_c+(1-t)\pi_c I$, with the preparation and commands
unchanged.

\begin{theorem}[Attainment of the module bound at small contrast]
\label{thm:r74-low-contrast}
For every finite observable group $G$ there is an explicitly positive $t_0$
such that
\begin{equation}\label{eq:r74-small-contrast}
 \kappa_G(b_t)=\nu_G(V_b),\qquad
 m(b_t)=\min\{N,1+\nu_G(V_b)\}\qquad(0<t\le t_0).
\end{equation}
For every $t>0$, the observable group, module and residual cardinality agree
with those of $b$. At $t=0$ the minimum is one label. The assertion allows
arbitrary finite groups and arbitrary finite terminal outcome sets.
\end{theorem}
\begin{proof}
Every $f\in V_b$ has a suffix-independent total mass
$\ell(f)=\sum_c f(z,c)$, because this holds for its spanning residuals.
Group averaging sends every residual to the same constant response
$\overline b(z)=\pi$. Its image on $V_b$ is therefore $\R\overline b$, and
\[
 V_b=\R\overline b\oplus W,
 \qquad W=\ker\left(\frac1{|G|}\sum_{g\in G}L_g\right).
\]
The module $W$ has no invariant vector, and $\ell(W)=0$.

Choose a $G$-set $X$ of minimum size $k=\nu_G(V_b)$ and an equivariant
injection $j:W\to\R^X$. The permutation inner product is invariant.
Since group averaging commutes with $j$, its image is orthogonal to the
whole invariant subspace, including the uniform row
$p=(1/k,\ldots,1/k)$. In particular its vectors have coordinate sum zero.
Let $F:\R^X\to W$ be orthogonal projection onto $j(W)$ followed by
$j^{-1}$. It is equivariant, $Fj=I_W$, and $F(p)=0$.

Put $v=b-\overline b$, $a=j(v)$, and
\[
 T(x,c)=F(e_x)(\varnothing,c).
\]
The rows of $T$ sum to zero, $pT=0$, and equivariance gives
$aP_gT=b(g)-\pi$. Since $a\ne0$, we may take
\[
 \epsilon=\frac{1}{2k\|a\|_\infty},\qquad
 B=\max_{x,c}\frac{|T(x,c)|}{\pi_c}>0,
 \qquad t_0=\min\{1,\epsilon/B\}.
\]
The row $\lambda=p+\epsilon a$ is a probability row with every entry at
least $1/(2k)$. Since $(t/\epsilon)B\le1$, for $0\le t\le t_0$ the matrix
\[
 D_t(x,c)=\pi_c+(t/\epsilon)T(x,c)
\]
is row-stochastic. Direct multiplication yields
\[
 \lambda P_gD_t
 =\pi+(t/\epsilon)\{pP_gT+\epsilon aP_gT\}
 =b_t(g).
\]
This proves the upper $\kappa_G(b_t)\le k$ with a bounded decoder.

The linear map
$f\mapsto tf+(1-t)\ell(f)\overline b$ is equivariant and invertible for
$t>0$: it is the identity on $\R\overline b$ and multiplication by $t$ on
$W$. It sends each residual of $b$ to the corresponding residual of $b_t$.
It preserves the group kernel and the equality pattern among residuals.
The module lower bound \eqref{eq:r73-module-bounds} now gives equality,
and Theorem~\ref{thm:r73-memory-formula} charges the pure start. At zero
contrast the decoder is constant.
\end{proof}

The construction is an equivariant positivity argument. Symmetric cone
factorizations and coordinate-symmetric lifts have an established theory
\cite{GPT13,KPT10}; ordinary nonnegative rank also has classical barycentric
perturbation results \cite{BocciCarliniRapallo11}. The assertion here is the
exact value $\nu_G(V_b)$ while preserving one common action for every command
word and the prescribed-start charge. It is not a claim that positivity
under mixing or equivariant factorization was previously unknown.

There is a useful distinction at an infinite compact observable group.
The same invertible change of residuals preserves that group for every
$t>0$, so Theorem~\ref{thm:r72-reversible} still excludes every finite
exact stochastic realization. Thus making the terminal probabilities
strictly interior need not make exact memory finite. Terminal readout
contrast, repeated channel damping, and perturbations of an entire
instrument are different operations.

\begin{theorem}[Attained exact contrast thresholds]
\label{thm:r74-contrast-thresholds}
For each integer $k\ge1$ there is $\tau_k\in[0,1]$ such that
\begin{equation}\label{eq:r74-threshold-set}
 \kappa_G(b_t)\le k\quad\Longleftrightarrow\quad 0\le t\le\tau_k.
\end{equation}
These thresholds are nondecreasing in $k$, satisfy $\tau_N=1$, and obey
\[
 k<\nu_G(V_b)\Longrightarrow\tau_k=0,
 \qquad \tau_{\nu_G(V_b)}>0.
\]
The lower memory value is retained at every threshold itself. There are
only finitely many positive-contrast transitions. If all $b(g,c)$ are real
algebraic, all thresholds are effectively computable real algebraic numbers.
When all subgroups of $G$ are normal, they are obtained by finitely many
linear programs; rational tables then have rational thresholds and rational
optimal controllers.
\end{theorem}
\begin{proof}
For a fixed action on at most $k$ points, the variables $\lambda,D,t$ lie
in a compact product of probability simplexes and $[0,1]$. The equations
$\lambda P_gD=b_t(g)$ define a closed set, so its feasible contrasts form
a compact set. If $s>0$ is feasible, the same $\lambda$ and the decoder
$(t/s)D+(1-t/s)\mathbf1\pi$ make every $0\le t\le s$ feasible.
There are only finitely many actions on at most $k$ labelled points;
their feasible sets therefore form a closed initial interval. A single
fixed point realizes $b_0$, including when no positive contrast is feasible.

The residual machine gives $\kappa_G(b_t)\le N$; the remaining inequalities
follow from Theorem~\ref{thm:r74-low-contrast} and its module lower bound.
The exact memory is recovered from
$m(b_t)=\min\{N,1+\kappa_G(b_t)\}$ for $t>0$.
For algebraic data the equations for each action are finite polynomial
constraints over the real algebraic numbers. Classical real quantifier
elimination computes their projected interval and endpoint, as in
Section~\ref{sec:r73-algebraic}. This is a proved finite procedure,
not a reported general quantifier-elimination execution. The final assertion
is the linear reduction and rationality proved in
Theorems~\ref{thm:r74-quotient-lp} and~\ref{thm:r74-quotient-dual} below.
\end{proof}

\subsection{A sharp transition for the retained prepared-qubit experiment}
The bounded-decoder obstruction of
Proposition~\ref{prop:r73-decoder-obstruction} has a complete contrast law.
Let $u=(1,-1,0)$ and
\[
 \beta_t(g)=\frac12+\frac t4uP_gu^{\mathsf T},\qquad g\in S_3,
\]
and let $b_t=(\beta_t,1-\beta_t)$. This is the same prepared-qubit
experiment with terminal effect $t|+\rangle\langle+|+(1-t)I/2$.

\begin{proposition}[An exact qubit readout threshold]
\label{prop:r74-s3-contrast}
For $0<t\le1$, this response has $d=3$, $N=6$, and
\[
 \kappa_{S_3}(b_t)=
 \begin{cases}3,&0<t\le2/3,\\6,&2/3<t\le1,\end{cases}
 \qquad
 m(b_t)=
 \begin{cases}4,&0<t\le2/3,\\6,&2/3<t\le1.\end{cases}
\]
At $t=0$ its memory is one.
\end{proposition}
\begin{proof}
The group, rank and residual count are unchanged for $t>0$ by
Theorem~\ref{thm:r74-low-contrast}. A faithful $S_3$-set of size at most
five has exactly one natural three-point orbit and otherwise only fixed
points or one two-point sign orbit, as in the proof of
Proposition~\ref{prop:r73-decoder-obstruction}. Let $A$ be the initial
mass on its natural orbit. Write its unnormalized initial row and scalar
decoder as
\[
 p=(A/3)\mathbf1^{\mathsf T}+r,
 \qquad D=q\mathbf1+s,
 \qquad r\mathbf1=0,\quad\mathbf1^{\mathsf T}s=0.
\]
The contribution is $Aq+rP_gs$. All remaining orbits contribute a linear
combination of the constant and sign functions. The functions $rP_gs$ and
$uP_gu^{\mathsf T}$ are orthogonal to both: their group sums vanish because
$|S_3|^{-1}\sum_gP_g=\mathbf1\mathbf1^{\mathsf T}/3$, and
$\sum_g\operatorname{sgn}(g)P_g=0$ entrywise. Consequently equality of the
full responses implies
\begin{equation}\label{eq:r74-s3-coefficient}
 rP_gs=(t/4)uP_gu^{\mathsf T}\qquad(g\in S_3).
\end{equation}

Here no unproved matrix-coefficient injectivity is needed. On the row plane
$\mathbf1^\perp$, use the basis $u,(1,1,-2)$. The matrices of the identity,
$(12)$, $(23)$ and the chronological cycle $(123)$ are respectively
\[
 I_2,\quad
 \begin{pmatrix}-1&0\\0&1\end{pmatrix},\quad
 \frac12\begin{pmatrix}1&1\\3&-1\end{pmatrix},\quad
 \frac12\begin{pmatrix}-1&1\\-3&-1\end{pmatrix}.
\]
They span all real two-by-two matrices: the first two span the diagonal
ones, and the latter two have independent off-diagonal parts. Therefore
\eqref{eq:r74-s3-coefficient} holds with every endomorphism of this plane
in place of $P_g$. Equality of the associated rank-one bilinear functionals,
with $t>0$, forces $r=au$, $s=bu^{\mathsf T}$ and $ab=t/4$.
Nonnegativity of $p$ gives $|a|\le A/3$, while the bounded decoder gives
$|b|\le\min(q,1-q)\le1/2$. Hence
\[
 t/4=|ab|\le A/6\le1/6.
\]
No core of size at most five is possible above $t=2/3$.

For $0<t\le2/3$ use the natural three-point action with
$\lambda=(2/3,0,1/3)$ and scalar decoder
$D=\tfrac12\mathbf1+(3t/4)u^{\mathsf T}$. Its entries lie in $[0,1]$
and direct multiplication gives $\beta_t$; the rank bound proves
$\kappa=3$. Above the threshold the six-residual machine supplies the
matching upper. The pure-start formula gives the stated minima and the
constant case is immediate.
\end{proof}

% Research draft for R75; deliberately outside the sealed R74 source.
% Depends on the retained permutation-core normal form and positivity argument.
\section{Orbit multiplicity, positive synthesis, and the convexity boundary}
\label{sec:r75-orbit-multiplicity}

The quotient reduction of Section~\ref{sec:r74-quotient-memory} has a precise
boundary. A single transitive orbit admits convexification without enlarging
the orbit exactly when its stabilizer is normal. For general finite groups,
irreducible multiplicities measure the obstruction and determine the exact
number of prescribed orbit copies at sufficiently small positive contrast.
All representations below are finite dimensional. We use classical complete
reducibility, matrix-coefficient orthogonality, and Frobenius reciprocity
\cite{Serre77}; their representation-theoretic content is not claimed as new.

Let $G$ be a finite group, not necessarily the faithful observable group of
the particular table under discussion. For a nonempty finite right $G$-set
$X$, let $\mathcal R_X(C)$ denote the probability tables
\[
 b(g)=\lambda P_gD,\qquad
 \lambda\in\Delta(X),\quad D(x,\cdot)\in\Delta(C),
\]
where $P_gP_h=P_{gh}$. Thus the same preparation and decoder are used for
every group element. Taking an observable quotient does not change the
table; keeping $G$ explicit lets us compare different allowed actions.
For this ambient-group formulation, $\kappa_G(b)$ denotes the minimum
$|X|$ over all such factorizations; an action need not be faithful unless
$G$ is the observable group.
Write $\widehat G$ for the complex irreducible representations, choose a
unitary representative $\rho$ of degree $d_\rho$, and define
\begin{equation}\label{eq:r75-fourier-block}
 \widehat b_c(\rho)=\frac1{|G|}\sum_{g\in G}b(g,c)\rho(g^{-1}),
 \qquad
 A_\rho(b)=\begin{pmatrix}
 \widehat b_{c_1}(\rho)\\ \vdots\\ \widehat b_{c_{|C|}}(\rho)
 \end{pmatrix},
 \qquad r_\rho(b)=\operatorname{rank}_{\mathbb C}A_\rho(b).
\end{equation}
Changing unitary bases does not change these ranks.

\begin{lemma}[Fourier multiplicity and real descent]
\label{lem:r75-fourier-multiplicity}
Let $V_b$ be the real span of the translates $b^h(g)=b(hg)$. The
multiplicity of the right irreducible $\rho$ in $V_b\otimes_{\mathbb R}
\mathbb C$ is $r_\rho(b)$. In particular $r_{\boldsymbol 1}(b)=1$ and
$r_\rho(b)\le d_\rho$. If $U$ and $V$ are real $G$-modules, then $U$
embeds equivariantly in $V$ if and only if every complex irreducible
multiplicity in $U\otimes\mathbb C$ is at most the corresponding one in
$V\otimes\mathbb C$.
\end{lemma}
\begin{proof}
The Fourier transform sends $b^h$ to the blocks
$A_\rho(b)\rho(h)$. Matrix-coefficient orthogonality identifies
$\mathbb C[G]$ with $\bigoplus_{\rho\in\widehat G}M_{d_\rho}(\mathbb C)$:
the group algebra can independently prescribe an arbitrary matrix on each
irreducible block. Consequently the $\rho$-part of the cyclic module
generated by $b$ identifies with
\[
 \{A_\rho(b)M:M\in M_{d_\rho}(\mathbb C)\}.
\]
If $A_\rho$ has rank $r$, a choice of a basis of its column space
identifies this right module with $r$ independent rows of length
$d_\rho$. It therefore consists of $r$ copies of $\rho$. For the trivial
representation the stacked block is the nonzero probability vector
$\pi=|G|^{-1}\sum_gb(g)$, so its rank is one.

Necessity of the second assertion follows by complexification. For
sufficiency, complete reducibility over $\mathbb C$ and the multiplicity
inequalities give a complex equivariant injection
$F:U\otimes\mathbb C\to V\otimes\mathbb C$. In real bases write
$F=A+iB$. Both $A$ and $B$ are real intertwiners. A maximal nonzero minor
of $F$ is the value at $z=i$ of a real polynomial minor of $A+zB$.
That polynomial is not identically zero, so some real $z$ makes this
minor nonzero. The real map $A+zB$ is the required injection. One may
choose $z$ among $0,1,\ldots,\dim U$, since the nonzero polynomial has
degree at most $\dim U$. This proves the assertion without excluding
complex or quaternionic representation types.
\end{proof}

For $H\le G$ put
\[
 m_\rho(H)=\dim_{\mathbb C}\rho^H
 =\operatorname{rank}\left(\frac1{|H|}\sum_{h\in H}\rho(h)\right).
\]
By Frobenius reciprocity these are the multiplicities of
$\mathbb C[H\backslash G]$. Let $q_H(b)$ be the least positive integer
$q$ for which $b$ belongs to $\mathcal R_{q(H\backslash G)}(C)$, where
$q(H\backslash G)$ is the disjoint union of $q$ copies; set $q_H(b)=\infty$
if no such factorization exists. This is an action-restricted resource.
It is not $\kappa_G(b)/[G:H]$ in general, since an unrestricted optimal
core may use other orbit types.

\begin{theorem}[Exact orbit multiplicity at small contrast]
\label{thm:r75-orbit-multiplicity}
Let $b:G\to\Delta(C)$, let
$\pi=|G|^{-1}\sum_gb(g)$, and set $b_t=\pi+t(b-\pi)$ for $0\le t\le1$.
For $H\le G$ define
\begin{equation}\label{eq:r75-copy-number}
 Q_H(b)=\max_{\rho\in\widehat G}
 \left\lceil\frac{r_\rho(b)}{m_\rho(H)}\right\rceil,
\end{equation}
with a positive numerator and zero denominator interpreted as $\infty$
and $0/0$ as zero. For every $t>0$,
$q_H(b_t)\ge Q_H(b)$. If $Q_H(b)<\infty$, an explicitly positive $t_0$
exists such that
\[
 q_H(b_t)=Q_H(b)\qquad(0<t\le t_0).
\]
If $Q_H(b)=\infty$, no positive contrast is realizable using any finite
number of copies of this orbit. At zero contrast one copy suffices.

More generally, choose one representative of each subgroup conjugacy
class. The unrestricted permutation-module bound has the exact finite
integer formulation
\begin{equation}\label{eq:r75-multiplicity-cover}
 \nu_G(V_b)=
 \min\left\{
  \sum_H q_H[G:H]:\ q_H\in\mathbb Z_{\ge0},\quad
  \sum_Hq_Hm_\rho(H)\ge r_\rho(b)\quad(\rho\in\widehat G)
 \right\}.
\end{equation}
Consequently this integer minimum equals $\kappa_G(b_t)$ for every
sufficiently small positive contrast, with the pure-start charge given
by Theorem~\ref{thm:r73-memory-formula} when $G$ is the observable group.
\end{theorem}
\begin{proof}
A factorization through a finite $G$-set $Y$ makes $V_{b_t}$ a quotient
of a submodule of $\mathbb R^Y$. Averaging an inner product over the
finite group makes every invariant subspace have an invariant orthogonal
complement; hence this subquotient embeds into $\mathbb R^Y$. The
contrast change is the identity on the invariant line generated by
$\pi$ and multiplication by $t$ on the zero-average part, so for $t>0$
its module is isomorphic to $V_b$. On $q(H\backslash G)$ the available
complex multiplicities are $qm_\rho(H)$, proving the lower bound and
the assertion when an active type is absent.

For the matching construction, let $Y$ be any nonempty finite $G$-set
whose complex multiplicities dominate $r_\rho(b)$. The lemma gives a
real equivariant injection of $V_b$ into $\mathbb R^Y$. Here are the
positivity details, to make clear that no complex factorization is
being substituted for a stochastic one. Delete the output columns
with $\pi_c=0$; nonnegativity makes them identically zero, and they can
be restored afterwards. Decompose
\[
 V_b=\mathbb R\overline b\oplus W,\qquad
 \overline b(g)=\pi,
\]
where $W$ is the kernel of group averaging. Every element of $W$ has
zero total output mass. Put $k=|Y|$, let $p$ be its uniform probability
row, and choose an equivariant injection $j:W\to\mathbb R^Y$.
The image of $j$ is orthogonal to all invariant vectors. Let
$F:\mathbb R^Y\to W$ be orthogonal projection onto $j(W)$ followed by
$j^{-1}$. Then $Fj=I_W$ and $F(p)=0$. Set
\[
 a=j(b-\overline b),\qquad T(y,c)=F(e_y)(e,c).
\]
Equivariance gives $aP_gT=b(g)-\pi$, while $pT=0$ and each row of $T$
sums to zero. If $b$ is nonconstant, choose
\[
 \epsilon=\frac1{2k\|a\|_\infty},\qquad
 B=\max_{y,c}\frac{|T(y,c)|}{\pi_c},\qquad
 t_0=\min\{1,\epsilon/B\}.
\]
The probability row $\lambda=p+\epsilon a$ and the row-stochastic
decoder $D_t(y,c)=\pi_c+(t/\epsilon)T(y,c)$ then satisfy
$\lambda P_gD_t=b_t(g)$ for $0\le t\le t_0$. For a constant table the
constant decoder works at every contrast. This is precisely the
retained positive construction, applied to the specified action
rather than only to a globally minimum one. Taking
$Y=Q_H(b)(H\backslash G)$ proves equality in the copy formula.

Every finite $G$-set is a disjoint union of coset orbits, and each such
orbit contributes the multiplicities $m_\rho(H)$. The real descent
lemma therefore makes the displayed integer constraints equivalent to
real module containment. Conversely every feasible integer vector
defines that disjoint-union action. This proves
\eqref{eq:r75-multiplicity-cover}; Theorem~\ref{thm:r74-low-contrast}
or the preceding construction gives its positive realization. The
regular orbit makes the minimum finite, since
$m_\rho(\{e\})=d_\rho\ge r_\rho(b)$.
\end{proof}

\begin{theorem}[Normality is the exact single-orbit convexity boundary]
\label{thm:r75-normality-convexity}
For $H\le G$ and the transitive right action $X=H\backslash G$, the
following conditions are equivalent:
\begin{enumerate}
\item $H$ is normal in $G$;
\item $\mathcal R_X(\{0,1\})$ is convex;
\item $\mathcal R_X(C)$ is convex for every finite outcome set $C$.
\end{enumerate}
If $H$ is normal, this response set is exactly the polytope of
probability tables constant on the left cosets of $H$. If $H$ is not
normal, nonconvexity occurs in every neighborhood of the interior
constant binary table $(1/d,1-1/d)$, where $d=[G:H]$.

An explicit witness is supplied by the return tables
\[
 f_x(g)=\mathbf1_{\{xg=x\}},\qquad x\in X,
 \qquad
 F_X(g)=\frac1d\sum_{x\in X}f_x(g)=\frac{\chi_X(g)}d.
\]
Every $(f_x,1-f_x)$ lies in $\mathcal R_X(\{0,1\})$, whereas
$(F_X,1-F_X)$ does not when $H$ is not normal. The same conclusion
holds after replacing every scalar $f$ by $d^{-1}+t(f-d^{-1})$, for
every $0<t\le1$.
More quantitatively, for any irreducible $\rho$ with
$0<m_\rho(H)<d_\rho$, put
$F_{X,t}=d^{-1}+t(F_X-d^{-1})$. Then
\begin{equation}\label{eq:r75-nonconvexity-margin}
 \inf_{(a,1-a)\in\mathcal R_X(\{0,1\})}
       \max_{g\in G}|F_{X,t}(g)-a(g)|
 \ \ge\ \frac{t\,m_\rho(H)}{d\,d_\rho}>0.
\end{equation}
This is a certified separation margin; optimality of its constant
is not asserted.
\end{theorem}
\begin{proof}
If $H$ is normal, for
$b(g,c)=\sum_{Hx}\lambda(Hx)D(Hxg,c)$ one has
$b(hg,c)=b(g,c)$ for every $h\in H$: normality gives $Hxh=Hx$.
Conversely every probability table with this left invariance is
realized by the point preparation $H$ and decoder
$D(Hg,c)=b(g,c)$. This proves the stated polytope identity and
(i)$\Rightarrow$(iii)$\Rightarrow$(ii).

Suppose $H$ is not normal. The group-algebra element
$e_H=|H|^{-1}\sum_{h\in H}h$ is an orthogonal projection in every
unitary irreducible representation, of rank $m_\rho(H)$. If every one
of these ranks were either zero or $d_\rho$, each matrix $\rho(e_H)$
would be zero or the identity. The faithful decomposition of the
complex group algebra into its irreducible matrix blocks would then
make $e_H$ central. Its coefficients are constant under conjugation
exactly when $H$ is normal. Hence some irreducible $\rho$ has
\[
 0<m_\rho(H)<d_\rho.
\]
This irreducible is nontrivial.

The return table at $x$ is realized by initial point $x$ and the
indicator decoder of $x$. Since
$\chi_X=\sum_\sigma m_\sigma(H)\chi_\sigma$, orthogonality gives
\[
 \widehat F_X(\rho)=\frac{m_\rho(H)}{d\,d_\rho}I_{d_\rho}.
\]
Its rank is $d_\rho>m_\rho(H)$. Lemma~\ref{lem:r75-fourier-multiplicity}
and the module lower bound therefore exclude a single-orbit
factorization. This average of members of the response set proves
nonconvexity. For positive contrast, the nontrivial Fourier block is
multiplied by $t$ and keeps its rank. Each contrasted return table is
still realized on its original orbit by replacing its decoder with
$d^{-1}+t(D-d^{-1})$. Letting $t$ decrease to zero proves the
neighborhood assertion. Thus (ii) implies (i).

For the quantitative assertion let $A=\widehat a(\rho)$ for any
single-orbit response. It has rank at most $m_\rho(H)<d_\rho$,
so there is a unit vector $v$ with $Av=0$. The Fourier block of
$F_{X,t}$ is $\alpha I$, where
$\alpha=t\,m_\rho(H)/(d\,d_\rho)$. Therefore
$\|\widehat F_{X,t}(\rho)-A\|_{\mathrm{op}}\ge\alpha$.
Because every $\rho(g)$ is unitary, the definition of the normalized
Fourier transform also gives
\[
 \|\widehat F_{X,t}(\rho)-\widehat a(\rho)\|_{\mathrm{op}}
 \le\frac1{|G|}\sum_{g\in G}|F_{X,t}(g)-a(g)|
 \le\max_{g\in G}|F_{X,t}(g)-a(g)|.
\]
Combining these inequalities proves the margin for every allowed
probability preparation and bounded decoder.
\end{proof}

For this witness the exact small-contrast number of copies is
\begin{equation}\label{eq:r75-character-copies}
 Q_H(F_X,1-F_X)
 =\max_{\rho:m_\rho(H)>0}
   \left\lceil\frac{d_\rho}{m_\rho(H)}\right\rceil.
\end{equation}
Indeed every active Fourier block of the character is full rank,
including the trivial block. Thus nonnormality entails a strict copy
cost at arbitrarily small positive contrast, with a matching positive
realization. This obstruction is not caused by a decoder touching the
boundary of its simplex.

\begin{lemma}[Full-contrast return-law rigidity]
\label{lem:r75-return-rigidity}
Let $X$ be a transitive $d$-point $G$-set. If a scalar response is a
convex combination of responses on copies of $X$, has value one at
$e$, and has group average $1/d$, every positive-weight component
uses a point preparation and its point indicator decoder. In
particular every component is a return law $f_x$.
\end{lemma}
\begin{proof}
For a component write $\lambda\in\Delta(X)$ and $D\in[0,1]^X$.
The total value one at $e$ forces $\lambda D=1$ for every component
of positive weight, hence $D(x)=1$ on the nonempty support of
$\lambda$. Transitivity gives its group average
$d^{-1}\sum_xD(x)\ge1/d$. Equality of the weighted mean with $1/d$
forces equality for each component. Consequently $\sum_xD(x)=1$,
so the support of $\lambda$ consists of exactly one point and $D$
is its indicator.
\end{proof}

These statements separate three matters: the classical multiplicities
of a permutation representation, the positivity of its preparation
and decoder, and the cost of an externally prescribed pure start.
The integer formula is exact at small positive contrast; it does not
replace the general full-contrast bilinear synthesis problem by a
linear program. The convexity theorem pinpoints why the retained
normal-stabilizer absorption cannot simply be extended to arbitrary
subgroups.

% Joint root/child research derivation for R75. Not a sealed source edit.
\subsection{An unbounded nonabelian family with a strict positivity cost}
\label{sec:r75-symmetric-return}

We apply the multiplicity theory to all symmetric groups $S_n$ with
$n\ge5$. The natural permutation character $\chi_n(g)$ counts the
fixed points of $g$ on $[n]$. Consider the binary response whose first
coordinate is
\begin{equation}\label{eq:r75-symmetric-return}
 \beta_t(g)=\frac{1-t}{n}+\frac{t}{n}\chi_n(g),
 \qquad b_t=(\beta_t,1-\beta_t),\qquad0\le t\le1.
\end{equation}
Its group average is $1/n$, and at the identity it equals
$[1+(n-1)t]/n$. This supplies a family with growing group, observable
rank, and exact memory, rather than a fixed-dimensional illustration.

The positivity threshold is controlled by the following matrix
extremum. No symmetry or positive-semidefiniteness hypothesis is
imposed on its coefficient matrix.

\begin{lemma}[A sharp singular-coefficient bound and its equality case]
\label{lem:r75-singular-coefficient}
Let $n\ge5$, $P=I-\mathbf1\mathbf1^{\mathsf T}/n$, and $t>0$.
If an entrywise nonnegative real $n\times n$ matrix $A$ satisfies
\[
 \operatorname{rank}A<n,\qquad
 \sum_{i,j}A_{ij}=1,\qquad PAP=\frac tnP,
\]
then
\begin{equation}\label{eq:r75-sharp-threshold}
 t\le\tau_n:=\frac{\sqrt n-1}{\sqrt n+1}.
\end{equation}
Equality holds if and only if, after the same permutation of rows
and columns,
\begin{equation}\label{eq:r75-extremal-coefficient}
 A=A_*:=\frac1{n-1}\sum_{i=1}^{n-1}
                  \ell_i^{\mathsf T}\ell_i,
 \qquad
 \ell_i=\left(1-\frac1{\sqrt n}\right)e_i+
                         \frac1{\sqrt n}e_n.
\end{equation}
Here $e_i$ and $\ell_i$ are row vectors. In particular symmetry of
the extremal matrix is a conclusion.
\end{lemma}
\begin{proof}
Put $\alpha=t/n$. The centered identity implies
$A=\alpha I+x\mathbf1^{\mathsf T}+\mathbf1y^{\mathsf T}$:
indeed $P(A-\alpha I)P=0$ decomposes the difference into a matrix
constant on each row and one constant on each column. Let $v$ be a
nonzero left null vector of $A$. If $v^{\mathsf T}\mathbf1=0$,
then $v^{\mathsf T}AP=v^{\mathsf T}PAP=\alpha v^{\mathsf T}$,
a contradiction. Normalize $\sum_i v_i=1$. Substitution of
$v^{\mathsf T}A=0$ in the preceding decomposition gives a vector
$c$ such that
\begin{equation}\label{eq:r75-null-vector-form}
 A=\alpha\{I+c\mathbf1^{\mathsf T}-\mathbf1v^{\mathsf T}\},
 \qquad v\cdot c=0,
 \qquad 1=t\sum_i c_i.
\end{equation}
We bound $\sum_i c_i$. Relabel so that $v_1=M$ is maximal and
$v_2=L$ is second largest. Since $\sum_i v_i=1$, $M>0$.
Entrywise nonnegativity gives
\[
 c_1\ge\max\{L,M-1\},\qquad c_i\ge M\quad(i>1).
\]
If $v$ is nonnegative, then $M\le1$ and $L\ge0$.
Thus $v\cdot c=0$ forces $v_i=0$ for every $i>1$,
$v_1=1$, and $c_1=0$. Consequently $\sum_i c_i\ge n-1$.

Otherwise put $m=-\min_i v_i>0$ and
$b_0=\max\{L,M-1\}$. Write
$c=(b_0,M,\ldots,M)+u$ with $u\ge0$. The dot product of $v$
with the displayed base vector is $M(b_0+1-M)\ge0$.
Since $v\cdot u\ge-m\sum_i u_i$, the condition $v\cdot c=0$
therefore yields
\begin{equation}\label{eq:r75-null-sum-bound}
 \sum_i c_i\ge b_0+(n-1)M+
                       \frac{M(b_0+1-M)}m.
\end{equation}
The normalization of $v$ also implies
\[
 M\ge M_-:=\frac{1+m}{n-1},\qquad
 L\ge\frac{1-M+m}{n-2}.
\]
The latter follows by removing one maximum and one minimum from
the sum of the coordinates and bounding each remaining coordinate
by $L$. The bounds $M-1$ and $(1-M+m)/(n-2)$ coincide at
$M_+:=1+m/(n-1)$.

If $M\ge M_+$, use $b_0\ge M-1$ in
\eqref{eq:r75-null-sum-bound} to obtain
$\sum_i c_i\ge nM-1\ge n-1$. If $M_-\le M\le M_+$,
use $b_0\ge(1-M+m)/(n-2)$. The right side of
\eqref{eq:r75-null-sum-bound} is increasing in $b_0$, and becomes
\[
 F(M)=\frac{1+m}{n-2}+(n-1)M+
              \frac{n-1}{m(n-2)}(M-M^2).
\]
This is a strictly concave quadratic. Its minimum on the stated
interval occurs at an endpoint, where
\[
 F(M_-)=\frac{1+m}{n-1}\left(n+\frac1m\right),
 \qquad F(M_+)=nM_+-1\ge n-1.
\]
The arithmetic--geometric mean inequality gives
\[
 \frac{1+m}{n-1}\left(n+\frac1m\right)
 =\frac{n+1+nm+1/m}{n-1}
 \ge\frac{(\sqrt n+1)^2}{n-1}.
\]
For $n>4$ this last constant is strictly less than $n-1$.
The two cases together prove
$\sum_i c_i\ge(\sqrt n+1)/(\sqrt n-1)$;
\eqref{eq:r75-null-vector-form} proves \eqref{eq:r75-sharp-threshold}.

Suppose equality holds, and write $s=\sqrt n$.
The case $v\ge0$ and the branch $M\ge M_+$ give strictly larger
bounds. Equality in the remaining chain forces
$m=1/s$ and $M=M_-$: strict concavity excludes the interior, and
the other endpoint has value at least $n-1$. Thus
\[
 M=\frac1{s(s-1)},\qquad
 v=(M,\ldots,M,-1/s)
\]
after a coordinate permutation. Indeed the sum of the $n-1$
coordinates other than one minimum is $(n-1)M$, forcing each
to equal its maximum. Now $b_0=M$. Equality in
$v\cdot u\ge-m\sum_i u_i$ places all of $u$ at the unique
negative coordinate, and its value there is $M/m$.
Hence $c=(M,\ldots,M,M+M/m)$. Substituting in
\eqref{eq:r75-null-vector-form} gives precisely
\eqref{eq:r75-extremal-coefficient}.

Conversely the rows $\ell_i$ in that formula are probability
vectors and linearly independent. If $p=\mathbf1^{\mathsf T}/n$
is regarded as a row, their centered rows satisfy
\[
 (\ell_i-p)\cdot(\ell_j-p)
    =\left(1-\frac1{\sqrt n}\right)^2\delta_{ij}.
\]
They are therefore an orthogonal basis of the zero-sum subspace.
Consequently $A_*$ is nonnegative, has total sum one and rank
$n-1$, and
\[
 PA_*P=\frac{(1-1/\sqrt n)^2}{n-1}P
       =\frac{\tau_n}{n}P.
\]
This proves sufficiency and completes the equality characterization.
\end{proof}

\begin{theorem}[The exact two-level nonabelian memory law]
\label{thm:r75-symmetric-return}
For $n\ge5$ and $t>0$, the response \eqref{eq:r75-symmetric-return}
has observable group $S_n$, residual cardinality $n!$, and observable
dimension $1+(n-1)^2$. Its permutation-module bound is
\[
 \nu_{S_n}(V_{b_t})=n(n-1).
\]
At the sharp threshold $\tau_n=(\sqrt n-1)/(\sqrt n+1)$, the
unrestricted positive core and pure-start memory satisfy
\[
 \begin{array}{c|c|c}
 \text{contrast}&\kappa_{S_n}(b_t)&m(b_t)\\\hline
 0<t\le\tau_n&n(n-1)&n(n-1)+1\\
 \tau_n<t\le1&n^2&n^2+1.
 \end{array}
\]
The smaller memory is attained at the threshold itself. At $t=0$
both minima are one.

Every minimum permutation core at $t=\tau_n$ is the disjoint union
of $n-1$ natural orbits. Its probability preparation and binary
decoder, after relabeling the orbit copies and simultaneously
relabeling their natural coordinates, have exactly the form
\begin{equation}\label{eq:r75-threshold-factorizations}
 \lambda_i=\ell_i,\qquad
 D_i=\frac{\ell_i}{(n-1)w_i},\qquad
 \sum_{i=1}^{n-1}w_i=1,\qquad
 w_i\ge\frac{1-1/\sqrt n}{n-1},
\end{equation}
where $\ell_i$ is defined in \eqref{eq:r75-extremal-coefficient},
$w_i$ is the preparation mass on orbit $i$, and $D_i$ is its
first-output decoder written as a row. The second-output decoder
is $1-D_i$.
Conversely every displayed weight vector gives an optimal
factorization. This is a classification of minimum permutation-core
factorizations under the inherited normal form, not a classification
of all stochastic transition realizations before that reduction.
\end{theorem}
\begin{proof}
The real natural permutation module is the direct sum of the
invariant line and the $(n-1)$-dimensional standard representation
$\mathrm{Std}_n$. Its complexification is irreducible on that
zero-sum subspace. One direct proof uses the transpositions: from
a nonconstant vector in an invariant subspace, a vector
$v-(ij)v$ produces a nonzero multiple of $e_i-e_j$, whose orbit
spans the whole zero-sum space. The standard representation is
faithful, because a permutation fixing every difference $e_i-e_j$
fixes all points. Matrix-coefficient orthogonality gives the
nontrivial block
\[
 \widehat\beta_t(\mathrm{Std}_n)
   =\frac{t}{n(n-1)}I_{n-1},
\]
and all other nontrivial irreducible blocks vanish. Therefore the
observable module is
$\boldsymbol1\oplus\mathrm{Std}_n^{\oplus(n-1)}$, of dimension
$1+(n-1)^2$. Equality of two residuals forces equality of their
standard representation matrices because this Fourier block is
invertible; faithfulness then makes the group elements equal.
This proves the claims about the group and residual count.

We next record the subgroup index estimate needed for an unrestricted
memory converse. For $K\le S_n$, let $r$ be the number of its orbits
on $[n]$. The multiplicity of $\mathrm{Std}_n$ in
$\mathbb C[K\backslash S_n]$ is
\[
 m(K)=\dim\mathrm{Std}_n^K=r-1.
\]
If $m(K)>0$, then
\begin{equation}\label{eq:r75-index-estimate}
 [S_n:K]\ge n\,m(K).
\end{equation}
More sharply, if $m(K)>0$ and $K$ is not a full point stabilizer,
\begin{equation}\label{eq:r75-index-extra}
 [S_n:K]\ge n\{m(K)+1\}.
\end{equation}
To prove these, let the natural orbit sizes of $K$ be
$a_1,\ldots,a_r$. Since $K$ is contained in the product of the
symmetric groups on those orbits,
\[
 [S_n:K]\ge\frac{n!}{a_1!\cdots a_r!}.
\]
For $r\ge3$, this is at least $n!/(n-r+1)!$: transferring one
unit from the smaller of two nonunit parts to the larger increases
their factorial product until only one nonunit part remains. The
last lower bound is at least $nr$ for $n\ge5$. Indeed at $r=3$
it is $n(n-1)\ge3n$, and increasing $r$ multiplies the bound by
at least two while $r$ increases by one. If $r=2$ and one part
is a singleton, $K$ lies in a point stabilizer; it is that full
stabilizer of index $n$, or has index at least $2n$. If both
parts have size at least two, the index is at least
$\binom n2\ge2n$. These cases prove both inequalities.
The argument uses the natural $K$-orbits and requires no classification
of low-index subgroups. In particular, an exceptional subgroup of
$S_6$ that is transitive in the natural action has $m(K)=0$ and
does not supply a standard carrier.

A $G$-set containing $V_{b_t}$ must have total standard multiplicity
at least $n-1$. Summing \eqref{eq:r75-index-estimate} over its
standard-carrying orbits gives at least $n(n-1)$ points.
Exactly $n-1$ natural orbits supply this multiplicity, so the module
minimum is $n(n-1)$.

For sharp positive attainment, use the preparations $\ell_i$ of
\eqref{eq:r75-extremal-coefficient} on $n-1$ natural orbits,
each of mass $1/(n-1)$, and use $D_i=\ell_i$ as first-output
decoders. The coefficient matrix is $A_*$. For any permutation
matrix $P_g$, the sum of its entrywise products with a matrix
$A$ will be denoted by $\langle A,P_g\rangle$; in particular
$\lambda P_gD^{\mathsf T}=\langle\lambda^{\mathsf T}D,P_g\rangle$
when the scalar decoder is written as a row.
The identities $PA_*P=(\tau_n/n)P$ and $\sum_{ij}(A_*)_{ij}=1$
imply
\[
 \langle A_*,P_g\rangle
 =\frac1n+\frac{\tau_n}{n}\{\chi_n(g)-1\}
 =\beta_{\tau_n}(g).
\]
Indeed the centered part contributes
$(\tau_n/n)(\chi_n(g)-1)$; the row-constant and column-constant
parts contribute their common group mean $1/n$.
Mixing each scalar decoder with $1/n$ by the factor $t/\tau_n$
realizes every $0\le t\le\tau_n$ on these same points.
The module bound proves optimality for positive $t$ in this interval.

A second, rational construction on the subinterval
$0\le t\le1/(n-1)$ is useful independently. Use $n-1$ natural orbits indexed by
$i=1,\ldots,n-1$, each with initial mass $1/(n-1)$ concentrated
at its point $i$. Give orbit $i$ the scalar decoder
\begin{equation}\label{eq:r75-symmetric-decoder}
 D_i(j)=\frac1n+\frac{(n-1)t}{n}
                 \{\mathbf1_{\{j=i\}}-\mathbf1_{\{j=n\}}\}.
\end{equation}
For $0\le t\le1/(n-1)$ these numbers lie in $[0,1]$.
Since a permutation has exactly one preimage of $n$,
\[
 \frac1{n-1}\sum_{i=1}^{n-1}D_i(g(i))
 =\frac1n+\frac tn
   \sum_{i=1}^{n-1}
    \{\mathbf1_{\{g(i)=i\}}-\mathbf1_{\{g(i)=n\}}\}
 =\frac1n+\frac tn\{\chi_n(g)-1\}=\beta_t(g).
\]
This retains a realization using only rational probabilities whenever
$t$ is rational. Its endpoint is not the sharp threshold when $n\ge5$.

At full contrast, $n$ natural orbits realize $\beta_1$: give
each mass $1/n$, start orbit $i$ at $i$, and use the indicator
decoder of $i$. Hence $\kappa_{S_n}(b_1)\le n^2$.
Suppose, to the contrary, that a factorization exists on fewer
than $n^2$ points. The total standard multiplicity is at least
$n-1$. If any standard-carrying orbit were not natural,
\eqref{eq:r75-index-extra} and the ordinary estimates on the
others would force at least $n(n-1)+n=n^2$ points. Thus all
standard-carrying orbits are natural. There must be exactly
$n-1$ of them: each supplies multiplicity one and costs $n$
points. Every other orbit has total size less than $n$.

Write the factorization as a mixture over its positive-mass
transitive orbits. Since $\beta_1(e)=1$, each component's
decoder equals one on its initial support. On a transitive orbit
of degree $d$ its mean is therefore at least $1/d$.
The natural components have mean at least $1/n$; any additional
orbit here has $d<n$ and would have mean strictly greater than
$1/n$. The mean of $\beta_1$ is exactly $1/n$, so all additional
orbits have zero initial mass and every active natural component
has mean $1/n$. Lemma~\ref{lem:r75-return-rigidity} now makes
each component a deterministic return law. Let $w_i$ be the
aggregate mass of the components starting at natural point $i$.
Evaluating at a transposition $(ij)$ gives
\[
 1-w_i-w_j=\frac{n-2}{n},
 \qquad w_i+w_j=\frac2n.
\]
For $n\ge3$ these equations imply $w_i=1/n$ for every $i$.
Thus at least $n$ positive components are needed, contradicting
the availability of only $n-1$ natural orbits. This proves
$\kappa_{S_n}(b_1)=n^2$.

We next exclude every intermediate positive-core value. Fix $t>0$
and suppose a factorization has fewer than $n^2$ points. The same
index argument gives exactly $n-1$ natural standard-carrying orbits;
the remaining orbits have standard multiplicity zero. Write
$b_t=B_N+B_E$ for their unnormalized response contributions. The
Fourier support of $B_N$ lies in the trivial and standard types,
because these are the only types in a natural permutation module.
The target has the same support. On the other hand, $B_E$ has zero
standard block. Subtracting $B_N$ from $b_t$ shows that all nontrivial
blocks of $B_E$ vanish, so $B_E(g)=v_E$ is a constant nonnegative
vector. Its total mass $w_E=\sum_c v_E(c)$ equals the preparation mass
on the extra orbits. Let $w_N=1-w_E$. Since $t>0$, the nonzero standard
block forces $w_N>0$. Restrict the original preparation and decoder to
the natural orbits, calling them $\lambda_N$ and $D_N$. On those same
$n(n-1)$ points use
\[
 \lambda'_N=\lambda_N/w_N,\qquad
 D'_N=w_ND_N+\mathbf1 v_E.
\]
Every row of $D'_N$ is nonnegative and sums to $w_N+w_E=1$, and
\[
 \lambda'_N P_gD'_N=B_N(g)+v_E=b_t(g).
\]
Thus any positive core below $n^2$ compresses to the module minimum
$n(n-1)$. No intermediate value is possible.

The full-contrast construction realizes every $b_t$ by mixing its
decoder with the constant $1/n$, so $\kappa\le n^2$ everywhere.
It remains to exclude the smaller core for every $t>\tau_n$.
By the absorption argument, any core below $n^2$ would give a
factorization on exactly $n-1$ natural orbits. Write its orbit
weights, probability rows, and scalar decoder rows as
$w_i$, $\lambda_i$, and $D_i$. The matrix
\[
 A=\sum_{i=1}^{n-1}w_i\lambda_i^{\mathsf T}D_i
\]
is entrywise nonnegative and has rank at most $n-1$.
Its response is $\langle A,P_g\rangle=\beta_t(g)$.
The group mean is $(1/n)\sum_{ij}A_{ij}=1/n$, so
$\sum_{ij}A_{ij}=1$. Moreover $PAP=(t/n)P$.
Here is an elementary way to see the last implication without a
choice of Fourier bases. Put $B=A-(t/n)I$. All permutation sums
$\langle B,P_g\rangle$ are equal. Comparing two permutations
that differ only by interchanging the assigned columns of two
distinct rows gives
\[
 B_{ij}+B_{k\ell}=B_{i\ell}+B_{kj}
 \quad(i\ne k,\ j\ne\ell).
\]
Thus $B_{ij}=B_{i1}+B_{1j}-B_{11}$, so $PBP=0$.
Lemma~\ref{lem:r75-singular-coefficient} now forces $t\le\tau_n$.
Together with the matching construction and the absence of
intermediate values, this proves the exact two-level law.
Theorem~\ref{thm:r73-memory-formula}, namely
$m=\min\{n!,1+\kappa\}$, gives the displayed pure-start values since
$n^2+1<n!$ for $n\ge5$. The explicit quadratic threshold is also
consistent with the retained algebraic-threshold theorem; no
quantifier-elimination calculation is needed here.

Finally consider a minimum core at $t=\tau_n$. The index bounds
force exactly $n-1$ natural orbits and no others, and the equality
part of Lemma~\ref{lem:r75-singular-coefficient} forces $A=A_*$
after simultaneous coordinate relabeling. Set
$a=1-1/\sqrt n$ and $b=1/\sqrt n$, so $a>b>0$ for $n\ge5$.
The leading $(n-1)\times(n-1)$ block of $A_*$ is the positive
diagonal matrix $a^2I/(n-1)$. Each nonnegative rank-one summand
$w_i\lambda_i^{\mathsf T}D_i$ can cover at most one of these
diagonal positions: covering two would give a positive off-diagonal
entry. All $n-1$ positions need coverage and there are exactly
$n-1$ summands. Therefore each is covered by a different
positive-weight summand. Relabel that summand by its position $i$.
The zero entries in row and column $i$ then force both
$\lambda_i$ and $D_i$ to be supported on $\{i,n\}$.
Only this summand contributes to the entries $(i,i)$, $(i,n)$,
and $(n,i)$. Their ratios in $A_*$ are
\[
 \frac{(A_*)_{in}}{(A_*)_{ii}}
 =\frac{(A_*)_{ni}}{(A_*)_{ii}}=\frac ba.
\]
Normalization of $\lambda_i$ and the diagonal value
$(A_*)_{ii}=a^2/(n-1)$ give
$\lambda_i=ae_i+be_n=\ell_i$ and
$D_i=\ell_i/\{(n-1)w_i\}$. The condition $D_i\in[0,1]^n$
is exactly $w_i\ge a/(n-1)$ because $a>b$.
Conversely these weights make every decoder valid and give
$w_i\lambda_i^{\mathsf T}D_i=\ell_i^{\mathsf T}\ell_i/(n-1)$.
This proves precisely the claimed classification.
\end{proof}

\paragraph{Prepared quantum realization.}
The family has a finite-dimensional prepared unitary realization,
so it remains inside the experiment framework. On
$\mathbb C^n\otimes\mathbb C^n$ take
\[
 \rho=\frac1n\sum_{i=1}^n|ii\rangle\langle ii|,
 \qquad U_g=P_g^{\mathsf T}\otimes I,
 \qquad E=\sum_{i=1}^n|ii\rangle\langle ii|.
\]
Here $P_g$ is the row-transition permutation matrix, so
$U_g|x,y\rangle=|xg,y\rangle$, consistently with chronological
command multiplication. Then $\operatorname{tr}(E U_g\rho U_g^*)=\chi_n(g)/n$.
The binary effect $E_t=tE+(1-t)I/n$ produces
\eqref{eq:r75-symmetric-return}. Two command types suffice uniformly in
$n$: an $n$-cycle and a transposition of adjacent points generate $S_n$,
since conjugation by the cycle supplies the adjacent transpositions.
They give one common experiment valid for every word in that fixed-size
command alphabet, with a binary terminal outcome set.
The pure maximally entangled preparation
$n^{-1/2}\sum_i|ii\rangle$ gives the same terminal probabilities
with this effect, if a pure preparation is desired. The statement
is a physical realization of the response, not a claim that this
Hilbert-space dimension is minimal or that it establishes a
quantum-versus-classical separation under a common resource metric.

The gap $n$ between the positive core at full contrast and its
module minimum is unbounded. The relative ratio $n/(n-1)$ tends
to one; no growing relative separation is asserted. The exact
transition tends to one as $n$ grows. Its proof combines a general
orbit-multiplicity converse, a sharp positivity inequality with
equality rigidity, and explicit optimal implementation; none of
these counts requires the completion of the wider Theta pipeline.

\section{Quotient decompositions, linear duality and sharp transitions}
\label{sec:r74-quotient-memory}

For groups whose subgroups are normal, the probability optimization in the
normal form has an exact linear reduction. This includes every finite abelian
group and allows arbitrary probability tables and terminal outcome sets.
Neither a special phase choice nor a positive-amplitude trigonometric ansatz
is imposed. The reduction uses normal stabilizers; it is not an unrestricted
claim that equivariant nonnegative factorization is linear programming.

\begin{lemma}[Absorption of within-orbit preparation]
\label{lem:r74-normal-orbits}
Suppose the stabilizers in a finite $G$-set are normal. Every probability
factorization on it can be replaced, without increasing its size, by one
having at most one orbit $H\backslash G$ for each normal subgroup $H$ and
whose initial row is concentrated at the identity coset within each orbit.
Only the mixture between different orbits is required.
\end{lemma}
\begin{proof}
On an orbit of positive initial mass, normalize its local preparation and
write its response as
\[
 q(g,c)=\sum_{Hx\in H\backslash G}\lambda(Hx)D(Hxg,c).
\]
For $h\in H$, normality gives $Hxhg=Hxg$, so $q(hg,c)=q(g,c)$.
Thus $q$ itself is a probability decoder on $H\backslash G$. Replacing
the preparation by the identity coset and the decoder by $q$ leaves every
command-word response unchanged. Orbits with the same $H$ can be merged
by averaging their decoders with their initial masses as weights.
Delete zero-mass orbits. These operations weakly decrease the size.
\end{proof}

Assume henceforth, until the Fourier specialization, that every subgroup of
$G$ is normal. Such groups are called Dedekind groups. For a nonempty set
$\mathcal H$ of distinct subgroups, put
$c(\mathcal H)=\sum_{H\in\mathcal H}[G:H]$.

\begin{theorem}[Exact quotient linear programs]
\label{thm:r74-quotient-lp}
The invariant $\kappa_G(b)$ is the minimum $c(\mathcal H)$ for which there
are $p_H\ge0$ and functions
$f_H:H\backslash G\times C\to\R_{\ge0}$ satisfying
\begin{equation}\label{eq:r74-quotient-lp}
 \sum_cf_H(q,c)=p_H\quad(q\in H\backslash G),\qquad
 \sum_Hp_H=1,\qquad
 \sum_H f_H(Hg,c)=b(g,c)\quad(g\in G,c\in C).
\end{equation}
The residual cardinality is $N=|G|$, and
$m(b)=\min\{|G|,1+\kappa_G(b)\}$.
For a rational response table a minimum controller can be chosen rational.
More generally, if the table lies in a specified real algebraic number field,
the controller can be chosen in that same field.
\end{theorem}
\begin{proof}
Apply Lemma~\ref{lem:r74-normal-orbits} to a minimum probability
factorization. Its orbit contributions give \eqref{eq:r74-quotient-lp}.
Conversely, if $p_H>0$, put $D_H=f_H/p_H$, place mass $p_H$ at its identity
coset, and use the right quotient action. If $p_H=0$, all its entries vanish
and the orbit is discarded. This constructs the whole response with at
most $c(\mathcal H)$ points. The regular orbit $H=\{1\}$ is always
available, and there are finitely many subgroup families.

The stabilizer $H_0$ of the initial residual is normal. If $h\in H_0$
and $g\in G$, moving $h$ past $g$ using normality shows that $h$ fixes
the residual after $g$. It therefore fixes their span and is the identity
in the minimal observable group. Thus $H_0=\{1\}$ and orbit-stabilizer
gives $N=|G|$. The pure-start assertion follows from
Theorem~\ref{thm:r73-memory-formula}.

For each family the displayed feasible set is a bounded closed polyhedron:
$0\le f_H(q,c)\le p_H\le1$. A feasible vertex is obtained by solving a
linear system over the coefficient field. It is rational for rational data
and remains in the given real algebraic field otherwise. The construction
uses only these entries and division by positive $p_H$. If the residual
machine is smaller than the initializer construction, its decoder is the
given table and has the same arithmetic property.
\end{proof}

This theorem optimizes the common permutation action supplied by the inherited
normal form. The relation between symmetry and positive factorizations is
classical \cite{GPT13,KPT10}, and positivity can cost more than a factorization
inside the target linear space \cite{GillisGlineur12}. The additional step
here is the absorption of preparation on normal quotient orbits, which
removes the bilinear unknowns without changing the counted resource.

\subsection{An explicit payoff dual}
Let $\mathcal P_H$ be the polytope of all probability tables on
$H\backslash G$, pulled back to $G$. For a real table $y$ on $G\times C$,
use the unnormalized pairing $\langle y,b\rangle=\sum_{g,c}y(g,c)b(g,c)$
and set
\begin{equation}\label{eq:r74-dual-support}
 h_H(y)=\sum_{q\in H\backslash G}\max_{c\in C}\sum_{g\in q}y(g,c),
 \qquad h_{\mathcal H}(y)=\max_{H\in\mathcal H}h_H(y).
\end{equation}

\begin{theorem}[Payoff certificates and exact contrast endpoints]
\label{thm:r74-quotient-dual}
The family $\mathcal H$ is feasible in \eqref{eq:r74-quotient-lp} if and
only if
\begin{equation}\label{eq:r74-payoff-test}
 \langle y,b\rangle\le h_{\mathcal H}(y)
 \quad\hbox{for every real table }y.
\end{equation}
If infeasible, it has a strict rational witness $y$. A witness of margin
$\Delta=\langle y,b\rangle-h_{\mathcal H}(y)>0$ still excludes that family
for any table $\widetilde b$ with
$\|\widetilde b-b\|_\infty<\Delta/\|y\|_1$.

For the contrast line \eqref{eq:r74-contrast}, put $v=b-\pi$. Its endpoint
for this family is
\begin{equation}\label{eq:r74-threshold-dual}
 \tau_{\mathcal H}
 =\min\left\{1,
  \min_{\langle y,v\rangle=1}
       \bigl(h_{\mathcal H}(y)-\langle y,\pi\rangle\bigr)\right\}.
\end{equation}
Both minima are attained. Equivalently, maximize $t$ in $[0,1]$ subject to
\eqref{eq:r74-quotient-lp} with right side $\pi+tv$.
The endpoint for budget $k$ is the maximum of these linear-program values
over $c(\mathcal H)\le k$. Rational tables give rational endpoints and
optimal controllers; no polynomial bound on the number of families is claimed.
\end{theorem}
\begin{proof}
The feasible response set for a fixed family is exactly
$\mathcal P_{\mathcal H}=\operatorname{conv}\bigcup_H\mathcal P_H$.
Each $\mathcal P_H$ is a product of probability simplexes. Maximizing a
payoff independently on each coset gives its support function $h_H$;
taking a convex hull of the union gives $h_{\mathcal H}$.
Membership proves \eqref{eq:r74-payoff-test}. If $b$ is outside this compact
convex set, let $q$ be a closest point. The first variation along the
segment from $q$ to any point of the set gives
$\langle b-q,x-q\rangle\le0$. With $y=b-q$ this yields a strictly
positive gap of at least $\|b-q\|_2^2$. Approximate $y$ by a rational
table sufficiently closely to preserve that gap; the support function is
a finite maximum of linear functions and is continuous. The perturbation
assertion is the inequality
$|\langle y,\widetilde b-b\rangle|\le\|y\|_1
\|\widetilde b-b\|_\infty$.

The constant table $\pi$ belongs to every $\mathcal P_H$. Along its ray
$\pi+tv$, a support inequality with $\langle y,v\rangle\le0$ imposes
no restriction for $t\ge0$; every other inequality can be normalized by
$\langle y,v\rangle=1$. This gives \eqref{eq:r74-threshold-dual}.
Attainment requires no strict-feasibility assumption. If $v$ lies outside
the linear space parallel to the affine hull of $\mathcal P_{\mathcal H}$,
an affine-hull equation gives an attained inner minimum zero. Otherwise
the bounded polytope has a finite facet description in that affine hull;
the first boundary reached along the ray activates a facet with positive
directional coefficient. Its normalized supporting functional attains
the inner minimum. The polytope is bounded and $v\ne0$, so the ray endpoint
is finite.

One explicit linear formulation of the inner dual introduces $z_{Hq}$
and $s$ with
\[
 z_{Hq}\ge\sum_{g\in q}y(g,c)\quad(H,q,c),\qquad
 s\ge\sum_qz_{Hq}\quad(H\in\mathcal H),\qquad
 \langle y,v\rangle=1,
\]
and minimizes $s-\langle y,\pi\rangle$. The primal formulation is obtained
by replacing the right side of \eqref{eq:r74-quotient-lp} by $\pi+tv$.
It is bounded and contains $t=0$. Vertex arithmetic proves the endpoint
and controller assertions, and the finite maximum over admissible families
gives the budget endpoint. These are exact linear certificates in the
specified finite problem, not a formal proof-assistant verification.
\end{proof}

\subsection{Fourier support and a sharp positivity transition}
For an abelian $G$, let $S$ be the complex characters whose vector Fourier
coefficient in $b$ is nonzero. It contains the trivial character and is
closed under conjugation.

\begin{corollary}[The abelian module cover]
\label{cor:r74-abelian-cover}
The number $\nu_G(V_b)$ is the minimum
$\sum_{H\in\mathcal H}[G:H]$ over distinct subgroup families such that
each character in $S$ is trivial on at least one $H\in\mathcal H$.
This lower bound is the exact value of $\kappa_G(b_t)$ for all sufficiently
small positive contrasts. It need not be attained at full contrast.
\end{corollary}
\begin{proof}
Character averaging projects the cyclic vector $b$ onto each of its
nonzero Fourier modes. Thus its complex cyclic module contains each
active character once; conjugate pairs give the corresponding real
two-dimensional summands. The permutation module on $H\backslash G$
contains exactly the characters trivial on $H$, each once. Containment
of $V_b$ is therefore exactly the stated cover. Repeated subgroups add no
new characters and may be deleted. Apply
Theorem~\ref{thm:r74-low-contrast}.
\end{proof}

The following example computes a complete rational threshold and an explicit
dual obstruction. Let $G=(\mathbb F_2)^4$, with the four coordinate flips as
commands, and let $\chi_1,\ldots,\chi_4$ be its coordinate sign characters.
Put $\chi_5=\chi_1\chi_2\chi_3\chi_4$ and
\begin{equation}\label{eq:r74-five-character-response}
 \beta_t(g)=\frac38+\frac t8\sum_{i=1}^5\chi_i(g),\qquad
 b_t(g)=(\beta_t(g),1-\beta_t(g)),\qquad 0\le t\le1.
\end{equation}

\begin{proposition}[A sharp transition invisible to the observable module]
\label{prop:r74-circuit-transition}
For every $t>0$, \eqref{eq:r74-five-character-response} has observable
group $(\mathbb F_2)^4$, dimension six, residual cardinality sixteen, and
$\nu_G(V_b)=10$. Its exact positive core and pure-start memory are
\[
 \kappa_G(b_t)=
 \begin{cases}10,&0<t\le3/5,\\12,&3/5<t\le1,\end{cases}
 \qquad
 m(b_t)=
 \begin{cases}11,&0<t\le3/5,\\13,&3/5<t\le1.\end{cases}
\]
At zero contrast the minimum is one. In particular a twelve-label minimum
does not occur in this family.
\end{proposition}
\begin{proof}
The five signs have product one, so their sum is $5,1$, or $-3$.
Consequently $\beta_1$ takes the values $1,1/2,0$, and every $\beta_t$
is a valid convex combination with the constant $3/8$.
Any four of the five characters are independent over $\mathbb F_2$.
For $t>0$ all five character coefficients and the constant mode are
nonzero, giving dimension six and trivial joint kernel. Thus $G$ is the
observable group and $N=16$.

For a quotient orbit write $W=H^\perp$ and $r=\dim W$; its size is $2^r$.
A proper $W$ contains at most $r$ of the five active nontrivial characters,
and $2^r\ge2r$ for $r=1,2,3$. An orbit with $r=4$ costs sixteen.
Covering the five modes therefore costs at least ten, and the five
two-point quotients attain this module bound.

Suppose that a positive core has at most eleven points. A size-eight
orbit covers at most three active characters; the remaining budget at
most three covers at most one more. Size sixteen is unavailable.
The nontrivial orbits thus have sizes two or four. Their total cost
must be ten and their active-character blocks partition the five modes
into singletons or pairs: each dimension must carry an active character
and no active character can be covered twice. At most one fixed point
can remain. Products belonging to different pairs are distinct and
inactive, because equality would give a nontrivial relation involving
at most four of the five circuit characters.

For this family of blocks let $\mathcal B_2$ denote its pairs and define
\[
 z(g)=1-\sum_{i=1}^5\chi_i(g)
       +\sum_{\{i,j\}\in\mathcal B_2}\chi_i(g)\chi_j(g).
\]
Averaging $z$ over the stabilizer of a singleton orbit gives $1-\chi_i$;
over a pair orbit it gives $(1-\chi_i)(1-\chi_j)$; over a fixed orbit
it gives $1$. The absence of any relation of length at most four ensures
that no other term survives these averages. Each resulting function is
nonnegative. Hence any nonnegative quotient contributions satisfy
\[
 \frac1{16}\sum_gz(g)\beta(g)
 =\sum_H\frac1{16}\sum_gz(g)f_H(Hg,1)\ge0.
\]
For the target \eqref{eq:r74-five-character-response}, character
orthogonality gives the value $(3-5t)/8$. Thus every core of size at
most eleven requires $t\le3/5$. This is also an explicit instance of
Theorem~\ref{thm:r74-quotient-dual}: take
$y(g,1)=-z(g)/16$, $y(g,0)=0$. Its family support value is zero and its
strict margin above the threshold is $(5t-3)/8$.

For $0<t\le3/5$, use five two-point quotient orbits, putting initial
mass $1/5$ at the plus point of each. The first-coordinate decoder on
orbit $i$ is
\[
 D_i(g,1)=\frac38+\frac{5t}{8}\chi_i(g).
\]
Its values lie in $[0,1]$, and the weighted sum is $\beta_t$. This
attains the ten-point core.

For $t=1$, use the four-point quotient on $(\chi_1,\chi_2)$ and the
eight-point quotient on $(\chi_3,\chi_4,\chi_5)$, with initial mass
$1/2$ at the all-plus point of each. Their decoders are respectively
the indicator of two plus signs and the indicator of at least two plus
signs. For signs $a,b,c$, the latter is
$1/2+(a+b+c-abc)/4$. The weighted sum of the two decoders is
\[
 \frac{1+\chi_1+\chi_2+\chi_1\chi_2}{8}
 +\frac14+
   \frac{\chi_3+\chi_4+\chi_5-\chi_3\chi_4\chi_5}{8}
 =\beta_1,
\]
because $\chi_1\chi_2=\chi_3\chi_4\chi_5$. Replace each decoder $D$ by
$(1-t)3/8+tD$ for other contrasts. This proves the twelve-point upper
everywhere, matching the strict lower above $3/5$.
The pure-start formula now gives eleven and thirteen, since $N=16$.
\end{proof}

\begin{corollary}[A robust finite-table obstruction]
\label{cor:r74-circuit-robust}
Fix $3/5<t\le1$ and a binary probability table $\widetilde\beta$ on the
same group $G=(\mathbb F_2)^4$. Supply it under the same four generator
commands, so the whole response is $w\mapsto\widetilde\beta(g(w))$.
If
\begin{equation}\label{eq:r74-circuit-robust}
 \max_{g\in G}|\widetilde\beta(g)-\beta_t(g)|
 <\frac{5t-3}{18},
\end{equation}
its exact positive core has at least twelve points and its pure-start
memory is at least thirteen. This is a sufficient lower certificate on
finite group tables. It does not assert a sharp neighbourhood radius or a
matching thirteen-label upper for every perturbed table.
\end{corollary}
\begin{proof}
Write $\eta$ for the left side of \eqref{eq:r74-circuit-robust}.
Every normalized Fourier coefficient changes by at most $\eta$.
Since $(5t-3)/18<t/8$ for $t\le1$, all five prescribed coefficients
remain nonzero. Their joint kernel is trivial, so the perturbed table
still has observable group $G$ and sixteen residuals. An action of size
at most eleven must therefore cover those same five modes, and the
preceding proof gives its singleton-and-pair partition even if additional
Fourier coefficients are present.

There are zero, one, or two pairs in such a partition. Every permutation
of the five circuit characters extends to a linear automorphism of the
dual group, because any four are a basis and the fifth is their sum.
Thus the absolute sum of its witness $z$ depends only on the number of
pairs. Direct evaluation on the sixteen sign patterns gives the following
values and multiplicities:
\[
 \begin{array}{c|l|c}
 \text{pairs}&\text{value : multiplicity of }z
       &16^{-1}\sum_g|z(g)|\\ \hline
 0&-4:1,\ 0:10,\ 4:5&3/2\\
 1&-3:1,\ -1:6,\ 1:4,\ 3:2,\ 5:3&17/8\\
 2&-2:5,\ 0:4,\ 2:2,\ 4:4,\ 6:1&9/4
 \end{array}
\]
The normalized pairing with $\widetilde\beta$ is at most
\[
 \frac{3-5t}{8}+\frac94\eta<0.
\]
It cannot have the required nonnegative quotient decomposition. Hence
$\kappa_G(\widetilde b)\ge12$, and
$m(\widetilde b)=\min\{16,1+\kappa_G(\widetilde b)\}\ge13$.
The upper thirteen on the exact contrast line was supplied by a specific
construction and is not inferred for arbitrary perturbations.
\end{proof}

For an exact physical realization one may use the sixteen-dimensional
diagonal model: prepare the identity basis state, use regular permutation
unitaries for the four generators, and take the diagonal terminal effect
with entries $\beta_t(g)$. This is already classical; no quantum advantage
is asserted for this example. The general low-contrast theorem, the normal
quotient linear programs, and this exact transition have different scopes.

\section{The retained instrument realization theory and its boundary}
\label{sec:r74-instruments}

The group results concern terminal responses. For an outcome at every call,
one must instead preserve each command--outcome history and every subsequent
decision made from that visible history. This stronger interface is already
covered, under a sufficient uniform posterior condition, by the earlier
instrument theory of this paper. We recall its exact scope before proving
the new conditional-memory classification.

For fixed known completely positive branches $\mathcal I_{a,y}$ whose
command sums are trace preserving, the hypothesis is
\[
 \mathcal I_{a,y}(X)\succeq
   \frac\eta h\tr\bigl(\mathcal I_{a,y}(X)\bigr)I
 \quad(X\succeq0),\qquad \eta>0.
\]
It is a condition on every recorded branch, including the zero-branch
convention, and follows when depolarizing noise of weight $\eta$ is applied
after each branch. In every finite dimension $h\ge2$, and for $k$ prescribed
preparations, Theorem~\ref{thm:r68-instrument} gives at most
\[
 k+\left\lceil(4h/\eta)^{h-1}\right\rceil
\]
hidden labels, with at most $h^2$ successors per branch row. All commands use
the same physical dictionary. The complete proof in
Section~\ref{sec:interior-instruments} gives the operator identity for every
history; it preserves all adaptive stopped laws, including nontermination,
and nonnegative costs. The hidden label is not exposed to the observer.
Corollary~\ref{cor:r68-risk} consequently preserves the risk of every common
report-based observer, including constrained boundary budgets; it does not
identify hidden simulator size with optimal observer memory.

The same retained theory includes a terminating algebraic construction
(Proposition~\ref{prop:r68-algebraic}), interiority after a finite block with
the clock and intermediate reports counted (Theorem~\ref{thm:r68-block}),
and matching powers of $\eta^{-1}$ for universal state-positive
representation (Theorem~\ref{thm:r68-noise-exponent}). The latter lower bound
is for generators which are actual density matrices. It is not imported
as a lower bound against arbitrary nonnegative hidden-state representations.
The family in Section~\ref{sec:r74-mixing-obstruction} will have a lower
bound against that unrestricted class.

Invariant positive cones and their stochastic realizations are classical
\cite{Heller65,MonrasWinter14}. Okay's quantum simulation framework explicitly
includes invariant-polytope updates and adaptive composition
\cite{Okay26}. Pure-state nets and approximation of the density-matrix body
have an established convex-geometric theory \cite{AubrunSzarek17}; extension
complexity is a different resource \cite{Fawzi21}. These ingredients, and
the paper's already proved interior-instrument theorem, are not asserted
as new discoveries of this revision. Their role is to fix the sufficient
condition and the resource convention against which the new sharp examples
are compared. The next common-contraction construction gives an explicit
certificate in supplied linear-response coordinates; it is also an
application of positive-realization geometry, not a replacement for the
retained instrument theorem.

\section{The exact cost of initializing a stationary response}
\label{sec:r75-stationary-initialization}

The distinction between a freely chosen initial distribution and a
prescribed initial label is not peculiar to a periodic example.  There is
an exact general law.  Its proof is an elementary elimination in the
space of future response functions, within the classical positive
realization framework \cite{Heller65,MonrasWinter14,MonrasWinter16}.  It does not require
hidden states to lie in the minimal linear realization, or require their
unobserved transition matrix to be irreducible.

Let $A$ be a nonempty finite command alphabet, with a finite report set
$Y_a$ for each $a\in A$.  A finite classical realization consists of
nonnegative matrices $T_{a,y}$ such that
\[
 T_a:=\sum_{y\in Y_a}T_{a,y},\qquad T_a\one=\one,
\]
and an initial probability row $\mu$.  For a command--report word
$w=(a_1,y_1)\cdots(a_n,y_n)$ write
$T_w=T_{a_1,y_1}\cdots T_{a_n,y_n}$ and
$p(w)=\mu T_w\one$.  Optional terminal readouts have decoder
columns $d_{b,z}\ge0$, with $\sum_zd_{b,z}=\one$, and response
$p(w;b,z)=\mu T_wd_{b,z}$.  All following equalities include every specified
terminal readout when these are part of the experiment.

Call the prepared response \emph{predictively stationary under $\nu$}
if $\nu$ is a probability row on $A$ with every $\nu(a)>0$ and
\begin{equation}\label{eq:r75-predictive-stationarity}
 \sum_{a\in A}\nu(a)\sum_{y\in Y_a}p((a,y)w)=p(w)
 \qquad\text{for every future word }w.
\end{equation}
When terminal readouts are specified, additionally require
\begin{equation}\label{eq:r75-terminal-stationarity}
 \sum_{a\in A}\nu(a)\sum_{y\in Y_a}p((a,y)w;b,z)=p(w;b,z)
 \qquad\text{for all }w,b,z,
\end{equation}
including the empty word $w$.
This is an equality of observable responses; it does not assume
$\mu\sum_a\nu(a)T_a=\mu$ in a particular hidden realization.  For one
command and no terminal readout it is ordinary stationarity.  A response
has a one-label realization precisely when its report probabilities
factor as $p(w)=\prod_jp_{a_j}(y_j)$, and every terminal readout has a
fixed distribution independent of the preceding word.  For one command
this is the i.i.d. case.

If a stationary hidden row is desired, it can be chosen without adding
labels. Put $\overline T=\sum_a\nu(a)T_a$. Every $\mu\overline T^j$
has the same future response by the displayed stationarity identities.
A subsequential limit of the Ces\`aro averages
$H^{-1}\sum_{j=0}^{H-1}\mu\overline T^j$ is stationary for
$\overline T$ and retains that response. Thus the freely mixed minimum
below agrees with the convention requiring a stationary hidden row for
the averaged command, and with the usual stationary HMM convention when
there is only one command.

Let $n_{\mathrm{mix}}(p)$ be the minimum number of labels with an arbitrary
initial probability row, and let $n_{\mathrm{pure}}(p)$ be the minimum when
one prescribed initial label is counted.  The minima are $+\infty$ when
there is no finite realization.

\begin{theorem}[Stationary initialization law]
\label{thm:r75-stationary-initializer}
Suppose \eqref{eq:r75-predictive-stationarity} holds, together with
\eqref{eq:r75-terminal-stationarity} when terminal readouts are included.
If the response has
a one-label realization, both minima are one.  Otherwise
\begin{equation}\label{eq:r75-initialization-law}
 n_{\mathrm{pure}}(p)=n_{\mathrm{mix}}(p)+1,
 \qquad (+\infty)+1:=+\infty.
\end{equation}
The lower holds for all real nonnegative transition tables, including
reducible and nonminimal linear realizations.  The upper processes the
first actual command while entering the core; it neither discards a
physical call nor inserts an uncounted preliminary randomization.
\end{theorem}
\begin{proof}
Take any finite pure-start realization, with initial label $o$, and let
$f_i$ denote the complete future response from hidden label $i$.  Thus
$f_i(w)=e_i^{\mathsf T}T_w\one$, with the analogous values for terminal
readouts.  Put $\overline T=\sum_a\nu(a)T_a$.  Stationarity gives the
functional identity
\begin{equation}\label{eq:r75-initial-functional-identity}
 f_o=\sum_j\overline T_{oj}f_j.
\end{equation}
If $\overline T_{oo}=1$, nonnegativity and full support of $\nu$ imply
$(T_a)_{oo}=1$ for every $a$.  All report transitions from $o$ then return
to $o$.  Their probabilities, and the terminal readouts at $o$, constitute
a one-label realization.  In the case under consideration this is
excluded.  Hence
\[
 \beta_j=\frac{\overline T_{oj}}{1-\overline T_{oo}}\quad(j\ne o)
\]
is a probability row and $f_o=\sum_{j\ne o}\beta_jf_j$.

Delete label $o$.  On the remaining labels define
\begin{equation}\label{eq:r75-elimination-tables}
 \widehat T_{a,y}(i,j)
   =T_{a,y}(i,j)+T_{a,y}(i,o)\beta_j,
 \qquad i,j\ne o,
\end{equation}
and retain their terminal decoders.  These matrices are nonnegative and
their sums over $y$ are stochastic.  The recursion
$f_i((a,y)w)=\sum_jT_{a,y}(i,j)f_j(w)$, with the functional identity for
$f_o$ substituted on its right side, is exactly the recursion for
\eqref{eq:r75-elimination-tables}.  Induction on word length, starting with
the empty word and each terminal decoder, proves that every retained
$f_i$ is unchanged.  Initial row $\beta$ therefore realizes $f_o$ with
one fewer label.  This proves $n_{\mathrm{mix}}\le
n_{\mathrm{pure}}-1$ whenever a finite pure realization exists.

Conversely, given a $k$-label mixed-start realization $(\mu,T_{a,y})$,
append one label $o$.  Its report transition on command $a$ enters the
core with row $\mu T_{a,y}$ and never returns to $o$; the core retains its
original transitions.  A terminal readout at $o$ has row $\mu d_{b,z}$.
The first physical command is already included in this row, so induction
gives the original law at every word length, including zero.  This uses
$k+1$ labels and proves the reverse inequality.  A finite pure realization
is itself a finite mixed one, so the two finiteness assertions are
equivalent.  The one-label exception is immediate.
\end{proof}

\begin{corollary}[Stationary processes]
\label{cor:r75-stationary-process}
For every stationary finite-alphabet process that is not i.i.d., the
minimum pure-start edge-emitting hidden Markov realization has exactly
one more label than the minimum freely mixed realization, with the same
infinite-value convention.  In particular, the $q$ versus $q+1$ conclusion
of Theorem~\ref{thm:r74-depolarizing-obstruction} is an instance of a
general stationary principle, while its phase-dependent $q$ remains a
separate spectral conclusion.
\end{corollary}

Equality of these complete response functions also preserves the law
under every common randomized policy for commands, stopping and terminal
readout.  For each finite stopped history, multiply the equal response by
the same policy probabilities.  Summation gives all finite stopping
events, and decreasing continuation events give the nontermination
probability.  Consequently all nonnegative costs measurable from that
record have equal expectations, including infinite expectations.  These
are process-realization label counts.  If the external policy has its own
persistent memory, that memory is counted separately in a closed-loop
implementation; it is not supplied free by the theorem.

\section{Sharp conditional memory with damping and adaptive commands}
\label{sec:r75-noisy-conditional}

The one-dependent family of Section~\ref{sec:r74-mixing-obstruction}
also admits a sharp positive theory under conditional damping.  The
extension below realizes the complete observed instrument, with one
physical dictionary shared by every command.  The classical sector
inequality and the scalar thresholds of Section~\ref{sec:r69-spectral}
are inherited ingredients, not new spectral-region results.  Earlier
noise-dependent process lower bounds, including an inverse-square-root
worst-phase dependence, are due to Fanizza, Lumbreras and Winter
\cite[Theorem 2 and Section 4.2]{Fanizza24}.  Here the lower and upper meet
for an explicit uniformly positive, exactly depolarizing instrument, and
the upper works simultaneously under adaptive rotation commands.

For a qubit set
\[
 U_\theta=\begin{pmatrix}
 \cos(\theta/2)&-\sin(\theta/2)\\
 \sin(\theta/2)&\cos(\theta/2)
 \end{pmatrix},\qquad
 E=\begin{pmatrix}3/4&0\\0&1/4\end{pmatrix},\qquad
 \tau=|0\rangle\langle0|,\qquad
 \mathcal D(X)=\tr(X)I/2.
\]
For $0\le\alpha\le1$ define
\begin{equation}\label{eq:r75-noisy-instrument}
 \begin{split}
 \Psi^{\alpha,\theta}_0(X)
   &=\tfrac18[\alpha U_\theta XU_\theta^*
                    +(1-\alpha)\mathcal D(X)],\\
 \Psi^{\alpha,\theta}_1(X)&=\tfrac14\tr(EX)\tau,\\
 \Psi^{\alpha,\theta}_2&=\mathcal D-\Psi^{\alpha,\theta}_0
                                      -\Psi^{\alpha,\theta}_1.
 \end{split}
\end{equation}
The initial state is $I/2$.  An allowed command specifies $\theta$; the
report is $s\in\{0,1,2\}$.  A command may be selected from the visible
history.  A finite known command alphabet is enough for every finite
lower witness below.  The same construction also gives Borel transition
tables for the entire angle circle; this last formulation counts labels,
not the encoding or exact sampling cost of a real-valued command.

This family does not satisfy the retained uniform-interior hypothesis
on every posterior, even when $\alpha<1$: report $1$ always prepares the
pure state $\tau$.  Repeated reports $1$ keep that pure state, so passing
to a fixed block length does not restore that hypothesis.  The finite
construction below instead uses the common geometry of the damped
rotation, the reset vertex and the remaining branch.

\begin{lemma}[Physical margins and predictive rank]
\label{lem:r75-noisy-physical}
For every $\alpha\in[0,1]$ and every angle, the maps
\eqref{eq:r75-noisy-instrument} form a quantum instrument, their sum is
exactly $\mathcal D$, and their report probabilities lie in
$\{1/8\}$, $[1/16,3/16]$ and $[11/16,13/16]$, respectively.
The prepared response is predictively stationary under each unobserved
command.  At a fixed angle its output process is stationary and
one-dependent.  When $\alpha>0$ and $\sin\theta\ne0$, its word Hankel rank
is three, and
\begin{equation}\label{eq:r75-noisy-cylinder}
 p(1\,0^n\,1)=\frac{8^{-n}}{64}
           \left(1+\frac{\alpha^n}{2}\cos(n\theta)\right)
 \qquad(n\ge0).
\end{equation}
At $\alpha=0$ the complete command response has rank two and is independent
of the angles.  For all $\alpha$, it has no one-label realization.
\end{lemma}
\begin{proof}
The first two Choi matrices are positive.  Their operator norms are at
most $1/4$ and $3/16$, while $J(\mathcal D)=I_4/2$.  Hence
$J(\Psi^{\alpha,\theta}_2)\succeq I_4/16$.  The asserted trace bounds
are unchanged by damping.  Since $\mathcal D(I/2)=I/2$, averaging any
unobserved command fixes the prepared predictive response, including
arbitrary subsequent commands and terminal POVMs.  At a fixed angle,
summing an intervening report applies $\mathcal D$ and therefore
factorizes the probabilities of the blocks on either side, exactly as in
Theorem~\ref{thm:r74-depolarizing-obstruction}.

Report $1$ has initial probability $1/8$ and prepares $\tau$.  Conditional
on $n$ zero reports its Bloch vector is
$\alpha^n(\cos(n\theta),\sin(n\theta))$, giving
\eqref{eq:r75-noisy-cylinder}.  All reached subnormalized states are real
symmetric, so the Hankel rank is at most three.  For $\alpha>0$ and
nonreal phase, the three modes $1/8$, $\alpha e^{i\theta}/8$ and
$\alpha e^{-i\theta}/8$ are distinct with nonzero coefficients.  The
$3\times3$ submatrix with entries $p(1\,0^{i+j}\,1)$ is therefore a
product of two nonsingular Vandermonde matrices and a nonsingular
diagonal matrix.  This proves rank three.
At $\alpha=0$ every reached state is diagonal, giving rank at most two;
the $2\times2$ cylinder minor has determinant $2^{-19}>0$.
Finally $p(1)=1/8$ and $p(11)=3/128\ne p(1)^2$ for every parameter,
excluding a one-label realization.
\end{proof}

Write $N_{\mathrm{mix}}(\alpha,\theta)$ for the least label count
reproducing every report word at one fixed angle, and define
$N_{\mathrm{pure}}$ with a counted pure start.  The same conclusions below
hold if every terminal qubit POVM is included: the uppers realize the
subnormalized quantum states themselves and the lowers use report words
alone.
Write $\mathcal K_M$ for the set of complex eigenvalues of $M\times M$
row-stochastic matrices, as in Section~\ref{sec:r69-spectral}.

\begin{lemma}[Relative cylinder obstruction and polygon realization]
\label{lem:r75-full-instrument-bounds}
For $\alpha>0$ and $\sin\theta\ne0$, every $M$-label mixed-start
realization satisfies
\begin{equation}\label{eq:r75-cylinder-spectrum}
 \alpha e^{i\theta}\in\mathcal K_M.
\end{equation}
Conversely, let $P_q$ be the regular $q$-gon with vertices
$e^{2\pi i j/q}$, $q\ge3$, one of which represents $\tau$.  If
$\alpha e^{i\theta}\in P_q$, its $q$ physical pure states form one
dictionary invariant under all three normalized branches, and give a
$q$-label mixed-start realization of the complete instrument.  A single
dictionary works for any collection of angles satisfying this condition.
\end{lemma}
\begin{proof}
For a candidate HMM put $v=\mu T_1$ and $d=T_1\one$, so that the cylinder
in \eqref{eq:r75-noisy-cylinder} is $vT_0^nd$.  Restrict $T_0$ to labels
reachable from $v$ and from which $d$ is eventually positive, obtaining
$S$.  For every pair $i,j$ in this productive set, fixed positive entry
and exit paths give constants $c_{ij}>0$ and integers $a_i,b_j\ge0$ with
\[
 c_{ij}(S^n)_{ij}\le p(1\,0^{a_i+n+b_j}\,1)\le C_{ij}8^{-n}.
\]
Thus $\rho(S)\le1/8$.  Cayley--Hamilton and the three distinct modes in
\eqref{eq:r75-noisy-cylinder} force
$\alpha e^{i\theta}/8$ to be an eigenvalue of $S$.
An irreducible diagonal block $B$ carrying it has spectral radius
$0<R\le1/8$.  The Perron normalization proved before
Proposition~\ref{prop:r70-relative-spectrum} gives a stochastic matrix
$Q$ with eigenvalue $\alpha e^{i\theta}/(8R)$ and eigenvector $u$.
For a stationary row $\pi$ of $Q$, $\pi u=0$ because this eigenvalue is
nonreal.  Consequently
\[
 Q'=8R Q+(1-8R)\one\pi
\]
is stochastic and has eigenvalue $\alpha e^{i\theta}$.  Its size is at
most $M$; padding by self-loop labels proves
\eqref{eq:r75-cylinder-spectrum}.  This normalizes only the irreducible
block carrying the mode, and does not assume a positive Perron vector for
a reducible productive core.

Identify the real Bloch plane with $\mathbb C$, taking $\tau$ to $1$.
Multiplication by $\alpha e^{i\theta}$ maps $P_q$ into itself whenever
that point is in $P_q$, by its invariance under multiplication by a
$q$th root.  This handles branch zero.  Branch one resets to the vertex
$1$.  If $r\in P_q$ and $h=\tr(\Psi_1(\sigma_r))\in[1/16,3/16]$, the
normalized third-branch Bloch vector is
\[
 \frac{-\alpha R_\theta r/8-h(1,0)}{7/8-h}.
\]
Its norm is at most $5/11$, strictly below the inradius
$\cos(\pi/q)\ge1/2$.  Thus it too lies in $P_q$.
Choose barycentric coefficients for each of these normalized images and
multiply them by the branch probability.  The resulting nonnegative
matrices $T_s$ have stochastic sum and satisfy, for every dictionary state,
\[
 \Psi_s(\sigma_i)=\sum_jT_s(i,j)\sigma_j.
\]
Uniform initial mixture represents $I/2$.  Induction proves this identity
for every subnormalized word state; terminal effects and common adaptive
policies therefore have the same laws.  Triangulation of the polygon
provides a piecewise continuous, hence Borel, choice when the angle varies.
\end{proof}

\begin{theorem}[Sharp noisy-instrument memory]
\label{thm:r75-noisy-memory}
For $M\ge3$, $0<\theta\le2\pi/M$ and $0<\alpha\le1$,
\begin{equation}\label{eq:r75-first-sector}
 N_{\mathrm{mix}}(\alpha,\theta)\le M
 \quad\Longleftrightarrow\quad
 \alpha\le\frac1{\cos\theta+\sin\theta\tan(\pi/M)}.
\end{equation}
At equality the mixed minimum is exactly $M$, and the pure minimum is
$M+1$.  In every $M$-label mixed realization at this boundary, its
zero-report matrix is diagonally similar to
$\frac18[(1-t)I+tC_M]$ for some $0<t\le1$, up to cyclic relabelling and
orientation.  No analogous claim is made for its other report matrices.

For $0<\alpha<1$ put
\begin{equation}\label{eq:r75-universal-count}
 K(\alpha)=\max\left\{3,
               \left\lceil\frac\pi{\arccos\alpha}\right\rceil\right\}.
\end{equation}
One dictionary of $K(\alpha)$ labels realizes simultaneously every
rotation command in \eqref{eq:r75-noisy-instrument}, with a mixed start.
This count is minimal for the full angle command family, and the
corresponding pure minimum is $K(\alpha)+1$.  In particular
\[
 \max_{\sin\theta\ne0}N_{\mathrm{mix}}(\alpha,\theta)=K(\alpha),
 \qquad
 \max_{\sin\theta\ne0}N_{\mathrm{pure}}(\alpha,\theta)=K(\alpha)+1.
\]
The same common upper holds for each finite known angle alphabet, with
matching lower whenever that alphabet contains a phase attaining the
displayed maximum.

At $\alpha=0$, the complete command response has exactly two labels with
a mixed start and three with a pure start.  At $\alpha=1$, the full angle
family has no finite realization; for one nonreal angle its finite-order
and infinite-order alternatives are those of
Theorem~\ref{thm:r74-depolarizing-obstruction}.  Whenever a relevant
minimum is finite, pure initialization adds exactly one label.
\end{theorem}
\begin{proof}
The lower in \eqref{eq:r75-first-sector} follows from
\eqref{eq:r75-cylinder-spectrum} and Lemma~\ref{lem:r69-sector}.
For the upper, the radial boundary of the edge between $1$ and
$e^{2\pi i/M}$ is the right side of
\eqref{eq:r75-first-sector}; use the polygon realization.  At that
boundary strict monotonicity of $\cot(\pi/M)$ excludes every smaller
size at least three; the rank-three lower excludes sizes one and two.
For any realizing size $M$, the irreducible block in the preceding
proof must use all $M$ labels.  Applying the sector inequality before
the final radial contraction gives
\[
 \frac{\alpha}{8R}
  [\cos\theta+\sin\theta\tan(\pi/M)]\le1,
\]
so $R\ge1/8$, and hence $R=1/8$.  Equality in the inherited sector lemma
therefore fixes its Perron normalization to the lazy directed cycle.
The stationary initialization theorem supplies the pure count without
requiring the original diagonal similarity to be scalar.

The disk of radius $\alpha$ lies in the regular $K(\alpha)$-gon, including
the equality case $\alpha=\cos(\pi/K(\alpha))$.  The preceding lemma
therefore constructs one physical dictionary for every angle.  If
$K=3$, any nonreal phase has rank three and requires three labels.  If
$K>3$, take $\theta=\pi/(K-1)$.  The definition of the ceiling gives
$\alpha>\cos(\pi/(K-1))$, so the first-sector criterion excludes
$K-1$ labels for this single phase.  It also excludes every smaller size
for the full command family.  The pure upper appends one initializer,
and the single critical phase, together with
Theorem~\ref{thm:r75-stationary-initializer}, gives the matching lower.

At $\alpha=0$, use the two diagonal density matrices with Bloch
coordinates $1$ and $-3/11$.  Their initial probabilities are $3/14$ and
$11/14$.  In that order the report matrices are
\begin{equation}\label{eq:r75-zero-damping-tables}
 T_0=\begin{pmatrix}3/112&11/112\\3/112&11/112\end{pmatrix},\qquad
 T_1=\begin{pmatrix}3/16&0\\19/176&0\end{pmatrix},\qquad
 T_2=\begin{pmatrix}0&11/16\\7/88&121/176\end{pmatrix}.
\end{equation}
Their sum has both rows $(3/14,11/14)$, and direct substitution verifies
the subnormalized state identities.  The rank-two lower and the
stationary initialization law give the asserted exact counts.
Finally the undamped full family includes a non-root phase, so its
infinite lower follows from the retained theorem.  For every finite
angle alphabet the prepared predictive state is fixed by every
unobserved command, and $p(11)\ne p(1)^2$; thus the exact initialization
law applies at all parameters, not just at spectral boundaries.
\end{proof}

For a fixed phase the lower in Lemma~\ref{lem:r75-full-instrument-bounds}
and its regular-polygon upper also transfer the already proved arithmetic
envelope of Theorem~\ref{thm:r71-envelope}, with the same constants, to
the complete observed instrument.  Consequently its two logarithmic
exponents and its badly-approximable criterion transfer as well.  This is
a change from a scalar response to an entire instrument, not a new
number-theoretic assertion.  No equality between the full-instrument
minimum and the scalar spectral minimum is asserted outside the proved
sectors and boundary cases.

\section{The causal task-order transfer theorem}\label{sec:central}
Suppose that every terminating state-stopped competitor reaches at least $B_0\ge1$ live labels under $\theta_*$. Such a certificate must concern actual physical command compatibility, not printed action names. Define
\begin{equation}\label{eq:widths}
\begin{array}{c|cc}
 i&m_-(M)&m_+(M)\\\hline
 \cp&M&\lfloor M/D_0\rfloor\\
 \ev&M&\lfloor M/D_0\rfloor\\
 \st&M-B_0&\lfloor(M-B)/D_0\rfloor .
\end{array}
\end{equation}
The event upper also requires $M\ge B$. The trivial state-stopped lower $B_0=1$ is valid for a procedure initially waiting to make a call.

\begin{theorem}[Affine task-order transfer with a paid reference and common construction]\label{thm:transfer}
Assume all the prepared-kernel, full-report, task-order, effort, occupation, resource and common-construction hypotheses of Section~\ref{sec:data}. Fix $N,M$, put $b_-=N/\kappa_-$, $b_+=N/\kappa_+$, and suppose $b_-\ge1$ and $m_-(M)\ge1$. Set
\[
 t_N=(1-A_N)_+/u_-.
\]
If $t_N>0$, every one of the three risks satisfies
\begin{equation}\label{eq:central-lower}
 R_i^\Theta(N,M)\ge
 \bigl[l_-t_N\Q_{m_-(M)}(b_-/t_N)-E_N\bigr]_+.
\end{equation}
If $t_N=0$, the asserted lower is zero, with no evaluation of $\Q$ at infinity. An empty legal class still has risk $+\infty$.

Suppose also $m_+(M)\ge1$, $b_+\ge1$, and, for event stopping, $M\ge B$. If
\begin{equation}\label{eq:beta}
 \beta_N=\frac{l_+b_+}{1+b_+a_+}\ge1,
\end{equation}
one common Borel program with actual expected cost at most $N$ in every model attains
\begin{equation}\label{eq:central-upper}
 R_i^\Theta(N,M)\le
 \frac{u_+(1+b_+a_+)}{l_+}\Q_{m_+(M)}(\beta_N)+b_+e_+.
\end{equation}
It commits almost surely. If $\beta_N<1$, the always-accept first-attempt construction instead gives $R_i^\Theta(N,M)\le u_+\Q_{m_+(M)}(1)+e_+$ at cost at most $\kappa_+\le N$.

The lower is against every legal competing strategy and uses one actual witness. The upper is one parameter-independent strategy for all models. The two effort constants need not agree. Neither direction assumes that the comparison experiment belongs to $\Theta$ or that a normalized predictive update is uniformly contracting.
\end{theorem}

\begin{corollary}[Scale-uniform relative transfer]\label{cor:relative}
Under Theorem~\ref{thm:transfer}, suppose $A_N=E_N=a_+=e_+=0$. If $l_+b_+\ge1$, then
\begin{equation}\label{eq:relative}
 \frac{l_-}{u_-}\Q_{m_-(M)}(u_-b_-)
 \le R_i^\Theta(N,M)
 \le\frac{u_+}{l_+}\Q_{m_+(M)}(l_+b_+).
\end{equation}
The factors $l_\pm,u_\pm$ may be fixed nonunit constants, independent of budget and resolution. In particular no term proportional to $b_\pm$ times a fixed absolute calibration error is required in this relative-order class.
\end{corollary}
\begin{proof}
Substitute $t_N=1/u_-$ and $\beta_N=l_+b_+$ in the theorem.
\end{proof}

\begin{corollary}[Separate exact checkpoint and autonomous laws]\label{cor:exact}
Suppose the certificates are exact: $l_\pm=u_\pm=1$, all defects vanish, $\kappa_-=\kappa_+=\kappa$, and exports are the identity. Then, for $N\ge\kappa$,
\[
 R_\cp^\Theta(N,m)=\Q_m(N/\kappa).
\]
The constructed autonomous upper uses $B+m$ state-stopped labels, or $\max(B,m)$ event-stopped labels. If the mandatory live count is sharp, $B_0=B$, then
\begin{align}
 R_\st^\Theta(N,B+m)&=\Q_m(N/\kappa)\quad(m\ge1),\label{eq:typed-exact}\\
 R_\ev^\Theta(N,M)&=\Q_M(N/\kappa)\quad(M\ge B).\label{eq:event-exact}
\end{align}
The subtraction by $B$ belongs only to the state-stopped model. Nonidentity exports require the expanded dictionary in \eqref{eq:widths}, unless an optimal dictionary is invariant under all exports. These exact stopping distinctions are inherited from the complete companion and are reproved here to fix the resource semantics.
\end{corollary}
\begin{proof}
The same lower and upper profile coincide in each asserted interface. Identity exports give $D_0=1$. The packing and converse widths are proved in Lemma~\ref{lem:roles}. No lower on live states is used to forbid event-stopped reuse.
\end{proof}

The familiar absolute-defect theorem is also a specialization: take $l_\pm=u_\pm=1$ while retaining $A_N,E_N,a_+,e_+$. Thus relative and absolute calibration are composed in one statement. The chosen metric is the common scored task throughout; it is not changed between the converse and the construction.

\section{Proof of the central theorem}
\begin{lemma}[Profile attainment]\label{lem:attainment}
For every finite $m\ge1$ and $b\ge1$, the minimum in \eqref{eq:profile} is attained. The profile is nonincreasing in $m,b$ and left-continuous in $b>1$.
\end{lemma}
\begin{proof}
Write a dictionary as an ordered tuple in the compact set $H^m$, permitting repetitions. The feasible densities are a weak-star closed subset of the radius-$b$ ball of $L^\infty(\nu)$: positivity, the upper bound and the mass equality are each closed under pairing with $L^1(\nu)$. They are therefore weak-star compact. As the tuple varies, its squared-distance function varies uniformly on $H$, since all points belong to a bounded set. Along a convergent product net, split the objective difference into this uniform variation and the weak-star pairing against the fixed bounded measurable loss at the limit tuple. This proves continuity of the objective and attainment. Nets avoid an unwarranted sequential compactness assumption. Monotonicity follows from inclusion of feasible sets. For left continuity take an optimizer $h$ at $b>1$ and let $h_\lambda=(1-\lambda)h+\lambda$. It has unit mass and cap $b_\lambda=(1-\lambda)b+\lambda<b$. Its objective tends to the optimum as $\lambda\downarrow0$, while monotonicity bounds the optimum at $b_\lambda$ between the optimum at $b$ and this objective.
\end{proof}

\begin{lemma}[Actual-history comparison measure]\label{lem:comparison}
Fix any robustly admissible competitor of cost at most $N$, with its fixed terminal dictionary $C$. There are finite nonnegative measures $\xi_c$ on $X$, $c\in C$, such that $\xi=\sum_c\xi_c$ obeys
\begin{align}
 \xi&\le (\E_*A_\tau)\nu\le b_-\nu,\label{eq:measure-cap}\\
 1&\le u_-\xi(X)+A_N,\label{eq:measure-mass}\\
 \E_*\|Z-\widehat r\|^2
 &\ge l_-\sum_{c\in C}\int d_c\,d\xi_c-E_N.
 \label{eq:measure-risk}
\end{align}
The measure $\xi$ is not asserted to be the terminal law of any comparison execution.
\end{lemma}
\begin{proof}
Adjoin an independent uniform random variable to the current report on both sides. This is current-call randomness, not a seed retained between calls; it is sufficient to represent the competitor's randomized output and stopping decision. A finite-output kernel has such a representation by cumulative probabilities. At a fixed reached productive history $h$ and chosen command $a$, the post-report decision then gives an acceptance indicator $q_h(y)$ and an output $r_h(y)\in C$ on acceptance. Give $r_h$ an arbitrary fixed value on rejection. The notation $y$ now includes the current independent random variable. Integrating the original certificates over that variable preserves their constants. Dependence on the entire earlier history is permitted in these functions.

For each Borel $D\subset X$, set
\begin{equation}\label{eq:xi-definition}
 \xi_c(D)=\E_*\sum_{j\le\tau:\ j\ {
m productive}}
 \int\one_{\{\pi_{H_{j-1},a_j}(y)\in D\}}
 q_{H_{j-1}}(y)\one_{\{r_{H_{j-1}}(y)=c\}}
 P^0_{H_{j-1},a_j}(dy,H).
\end{equation}
The outer expectation is the actual pre-call occupation at $\theta_*$. All terms are jointly measurable and nonnegative. Tonelli defines the measures, whose total is finite because $\E_*A_\tau\le b_-$. Summing over $c$ and dropping the acceptance indicator gives \eqref{eq:measure-cap}: each comparison row has mark law $\nu$.

The analogous sum with the actual rows instead of $P^0$ has mass exactly one, because the actual procedure commits exactly once almost surely. Apply \eqref{eq:lower-mass} at each actual reached history, then sum. The occupation certificate yields \eqref{eq:measure-mass}. Apply \eqref{eq:lower-task} in the same way to the disjoint actual commitment events. Their expected losses sum to the actual risk. The comparison contribution is $d_{r_h(y)}(\pi_{h,a}y)$ by \eqref{eq:full-moments}; this identity conditions on the full current report and therefore does not overlook extra information seen by the competitor. The accumulated loss defect is at most $E_N$, giving \eqref{eq:measure-risk}. One may first truncate the number of calls. Monotone convergence for the nonnegative contributions and the finite defect bounds justify the limit.

No actual past is replaced by a comparison past. In particular the competitor need not terminate, or obey a cost bound, if hypothetically run under the comparison experiment. This distinction is why the argument applies when that experiment is not an operating model.
\end{proof}

\begin{remark}[Why the full report is required]
A constant mark can conceal a visible fair bit which also equals the future audit. A comparison using only that mark has $f=0,w=1$, but a two-value decoder reading the bit predicts it with zero loss. The hypotheses therefore condition on the full report seen by the competitor; matching a compressed mark marginal alone cannot establish the lower.
\end{remark}

\begin{lemma}[Live roles, output roles and physical compatibility]\label{lem:roles}
A program with $\ell$ live roles and $t$ output roles has a state-stopped implementation with $\ell+t$ labels and an event-stopped implementation with $\max(\ell,t)$ labels, preserving every path cost and complete stopped scored law. Conversely an event-stopped $M$-label program has at most $M$ decoder values. A state-stopped $M$-label program with at least $B_0$ reached live labels has at most $M-B_0$ terminal decoder values.

For a physical live-state certificate, consider finitely many mandatory pre-call contexts $I$ with no free context input. Suppose the commands legal at a history of context $i$ are exactly $\mathcal U_i$, and every terminating competitor visits every context with positive probability under the same witness. This positive probability may depend on the competitor; no uniform lower bound over competitors is required. Let $b_{\rm phys}$ be the least size of a cover of $I$ by subsets $J$ with $\bigcap_{i\in J}\mathcal U_i\ne\varnothing$. Then at least $b_{\rm phys}$ live labels are required. This certificate is invariant under duplicate names and bijective stateless recodings of the same physical commands.
\end{lemma}
\begin{proof}
For state stopping take disjoint copies of the live and terminal role sets and type them zero and one. For event stopping embed the live and terminal role sets separately into one alphabet of size $\max(\ell,t)$. On continuing transitions use the live embedding; on commitment emit the final event and use the terminal embedding. A command row is read only on a continuing branch, while the decoder is read only after commitment. Sharing their label therefore causes no conflict and uses no continuing mode bit. Couple the original and recoded coins. Every finite stopped prefix and every path cost agree, and almost-sure termination gives equality of the stopped scored laws. A fixed decoder has one value per retained label. In state stopping the global halt map excludes each reached live label from the output set, proving its subtraction. Event stopping has no such exclusion.

A live label reached in contexts $J$ must have one command distribution supported on all the corresponding $\mathcal U_i$. Since $J$ is finite, the union of the exceptional null command sets is null; hence their intersection supports a probability and is nonempty. The sets of contexts served by the labels cover $I$, proving the bound. Conversely a nonempty intersection contains a command and hence a Dirac command distribution. Replacing names without changing executed kernels, legal domains, costs and reports does not change these intersections. A translator which itself remembers a hidden phase is not such a recoding; its memory must be charged.
\end{proof}

\begin{proof}[Proof of Theorem~\ref{thm:transfer}]
Fix a competing program. Lemma~\ref{lem:roles} bounds its decoder width by $m_-(M)$. Let $s=\xi(X)$ in Lemma~\ref{lem:comparison}. If $t_N>0$, then \eqref{eq:measure-mass} gives $s\ge t_N$. Scale the measure to $\xi'=(t_N/s)\xi$. It has mass $t_N$ and remains bounded above by $b_-\nu$. Thus $h=(d\xi'/d\nu)/t_N$ has unit mass and cap $b_-/t_N$. Nearest-dictionary projection and nonnegativity imply
\[
 \sum_{c\in C}\int d_c\,d\xi_c
 \ge\int [w+\dist(f,C)^2]\,d\xi
 \ge t_N\Q_{m_-(M)}(b_-/t_N).
\]
Together with \eqref{eq:measure-risk} this proves the asserted lower for this competitor. Its worst-model risk dominates its risk at $\theta_*$. Taking the infimum over competitors gives \eqref{eq:central-lower}, without exchanging an infimum and a supremum. Nonnegative loss gives the zero truncation and the stated case $t_N=0$.

For the upper assume \eqref{eq:beta}. Choose an attained optimizer $(C,h)$ at width $m_+(M)$ and cap $\beta_N$. Let $r_C(x)$ be the first nearest point to $f(x)$ in a fixed ordering of $C$. After each attempt of $\mathcal G$, accept $x$ with probability $h(x)/\beta_N$, export $T_Jr_C(x)$ and commit; otherwise return to the start routing label. All these operations are Borel. The choice $J$ is used within the current attempt; any part retained thereafter is charged through the live role or exported decoder, not a free seed. The exports are assumed jointly measurable and take at most $D_0|C|$ distinct values, so their finite decoder and transitions are Borel as well.

Conditional on every surviving attempt history, the acceptance probability is at least
\begin{equation}\label{eq:hazard}
 p=\frac{l_+}{\beta_N}-a_+=\frac1{b_+}>0.
\end{equation}
Consequently the probability of reaching attempt $j+1$ is at most $(1-p)^j$, by induction on conditional survival probabilities. In particular commitment is almost sure and the expected number of attempts is at most $1/p=b_+$. Each begun attempt has conditional expected cost at most $\kappa_+$, not merely a cost conditional on a favorable outcome. Summing the predictable attempt costs gives actual expected cost at most $\kappa_+b_+=N$.

The conditional unnormalized loss contributed by an attempt is at most $u_+\Q_{m_+(M)}(\beta_N)/\beta_N+e_+$. Sum its contributions on the disjoint terminal events, using the same survival bound. The result is at most
\[
 b_+\left[\frac{u_+}{\beta_N}\Q_{m_+(M)}(\beta_N)+e_+\right],
\]
which is \eqref{eq:central-upper}. The program has $B$ live roles and at most $D_0m_+(M)$ terminal roles. Lemma~\ref{lem:roles} gives exactly the sufficient widths in \eqref{eq:widths}. A checkpoint counts only the terminal width. No physical renewal law was used: the conditional hazard controls survival even if the hidden state after rejection differs arbitrarily from its former state.

If $\beta_N<1$, commit after the first attempt, and use an optimizer at cap one. The loss certificate with $q=1$ gives $u_+\Q_{m_+(M)}(1)+e_+$, and the whole-attempt cost is at most $\kappa_+$. This proves all remaining assertions.
\end{proof}

\section{Scope and further mathematical questions}\label{sec:r73-scope}\label{sec:r74-scope}\label{sec:r75-scope}
The exact small-contrast formula applies to every finite observable group.
The normality theorem characterizes convexity for a fixed transitive
permutation action; it does not classify every possible linear reduction
of an arbitrary nonabelian realization problem. The symmetric-group family
has a closed optimal resource law at every contrast, including the
transition and its extremal permutation-core factorizations. A closed
high-contrast formula for arbitrary finite groups remains open.

The stationary initialization theorem applies to arbitrary finite-alphabet
stationary response laws and makes no quantum, reversibility, or finite-rank
assumption. The noisy-instrument formulas concern the displayed qubit
family. Its single controller works for every word over the supplied
command alphabet, including adaptive visible policies. One-dependence and
stationarity are assertions about the fixed-command process; adaptive
choice of commands can itself transmit dependence. The stopped-law
approximation is a uniform upper bound under a target expected-call budget,
not a matching lower bound for every approximation problem.

The retained general instrument theorem assumes uniform interiority of
normalized branches. It supplies a simulator of the experiment; its
private label is not a new observation available to the original policy.
Consequently neither that theorem nor the new simultaneous construction
classifies optimal feedback or tied boundary decisions for every causal
experiment. The general-$M$ squared-sign minimax-risk frontier and the
intermediate phase-type timing frontier remain distinct unresolved
problems. Program length, arithmetic precision, working storage and exact
sampling costs are also distinct from the number of live labels.

The full actual-history comparison theorem below and the independent
causal, reset, discrimination and geometric results retain their original
quantifiers. Local geometric conclusions remain local. Independent A1,
A2 and further Theta construction objectives require their own proofs;
none is inferred from the present positive-realization results. The
supplement contains complete technical arguments and the retained
independent mathematical developments. Historical drafts and proof
records have separate immutable locations and are not prerequisites for
reading the current theorem chain.

\begin{thebibliography}{99}
\bibitem{Blackwell} D. Blackwell, \emph{Equivalent comparisons of experiments}, Ann. Math. Statist. \textbf{24} (1953), 265--272. DOI: 10.1214/aoms/1177729032.
\bibitem{CGM97} J. A. Cuesta-Albertos, A. Gordaliza and C. Matr\'an, \emph{Trimmed $k$-means: an attempt to robustify quantizers}, Ann. Statist. \textbf{25} (1997), 553--576. DOI: 10.1214/aos/1031833664.
\bibitem{GG99} L. A. Garc\'ia-Escudero and A. Gordaliza, \emph{Robustness properties of $k$ means and trimmed $k$ means}, J. Amer. Statist. Assoc. \textbf{94} (1999), 956--969. DOI: 10.1080/01621459.1999.10474200.
\bibitem{SC} C. R. Shalizi and J. P. Crutchfield, \emph{Computational mechanics: pattern and prediction, structure and simplicity}, J. Statist. Phys. \textbf{104} (2001), 817--879. DOI: 10.1023/A:1010388907793.
\bibitem{BC} N. Barnett and J. P. Crutchfield, \emph{Computational mechanics of input--output processes: structured transformations and the $\varepsilon$-transducer}, J. Statist. Phys. \textbf{161} (2015), 404--451; arXiv:1412.2690.
\bibitem{AMTW} A. Ambainis, M. Mosca, A. Tapp and R. de Wolf, \emph{Private quantum channels}, Proc. 41st Annual Symposium on Foundations of Computer Science (2000), 547--553. DOI: 10.1109/SFCS.2000.892142.
\bibitem{Buscemi} F. Buscemi, \emph{Comparison of quantum statistical models: equivalent conditions for sufficiency}, Comm. Math. Phys. \textbf{310} (2012), 625--647. DOI: 10.1007/s00220-012-1421-3.
\bibitem{BR14} N. Bertschinger and J. Rauh, \emph{The Blackwell relation defines no lattice}, Proc. IEEE International Symposium on Information Theory (2014), 2479--2483. DOI: 10.1109/ISIT.2014.6875280; arXiv:1401.3146.
\bibitem{Banerjee05} A. Banerjee, X. Guo and H. Wang, \emph{On the optimality of conditional expectation as a Bregman predictor}, IEEE Trans. Inform. Theory \textbf{51} (2005), 2664--2669. DOI: 10.1109/TIT.2005.850145.
\bibitem{Stengel96} B. von Stengel, \emph{Efficient computation of behavior strategies}, Games Econom. Behav. \textbf{14} (1996), 220--246. DOI: 10.1006/game.1996.0050.
\bibitem{Balder89} E. J. Balder, \emph{On compactness of the space of policies in stochastic dynamic programming}, Stochastic Process. Appl. \textbf{32} (1989), 141--150. DOI: 10.1016/0304-4149(89)90058-6.
\bibitem{Barry14} J. Barry, D. T. Barry and S. Aaronson, \emph{Quantum partially observable Markov decision processes}, Phys. Rev. A \textbf{90} (2014), 032311. DOI: 10.1103/PhysRevA.90.032311; arXiv:1406.2858.
\bibitem{Ying18} S. Ying and M. Ying, \emph{Reachability analysis of quantum Markov decision processes}, Inform. Comput. \textbf{263} (2018), 31--51. DOI: 10.1016/j.ic.2018.09.001; arXiv:1406.6146.
\bibitem{Saldi25} N. Saldi, S. Sanjari and S. Y\"uksel, \emph{Quantum Markov decision processes: general theory, approximations, and classes of policies}, SIAM J. Control Optim. \textbf{63} (2025), 1156--1183. DOI: 10.1137/24M1643864; arXiv:2402.14649.
\bibitem{Seidenberg54} A. Seidenberg, \emph{A new decision method for elementary algebra}, Ann. of Math. (2) \textbf{60} (1954), 365--374. DOI: 10.2307/1969640.
\bibitem{Chatterjee15} K. Chatterjee, M. Chmel\'ik, R. Gupta and A. Kanodia, \emph{Optimal cost almost-sure reachability in POMDPs}, Proc. AAAI Conference on Artificial Intelligence \textbf{29} (2015). DOI: 10.1609/aaai.v29i1.9683.
\bibitem{Andriushchenko23} R. Andriushchenko, A. Bork, M. \v{C}e\v{s}ka, S. Junges, J.-P. Katoen and M. Mac\'ak, \emph{Search and explore: symbiotic policy synthesis in POMDPs}, Computer Aided Verification (2023); arXiv:2305.14149v2.
\bibitem{Halmos69} P. R. Halmos, \emph{Two subspaces}, Trans. Amer. Math. Soc. \textbf{144} (1969), 381--389. DOI: 10.1090/S0002-9947-1969-0251519-5.
\bibitem{Kessler21} D. A. Kessler, E. Barkai and K. Ziegler, \emph{First-detection time of a quantum state under random probing}, Phys. Rev. A \textbf{103} (2021), 022222. DOI: 10.1103/PhysRevA.103.022222; arXiv:2012.01763.
\bibitem{DelVecchio26} G. Del Vecchio Del Vecchio and S. N. Majumdar, \emph{Optimal detection of quantum states via projective measurements}, J. Phys. A: Math. Theor. \textbf{59} (2026), 035001. DOI: 10.1088/1751-8121/ae34fe; arXiv:2509.08556.
\bibitem{Companion} Q. Qi, \emph{General Theta Foundations I: Acquired Geometry and Causal Resource Transfer}, complete technical companion on absolute calibration, exact stopping resources and nonintervening sensing, 2026. The complete retained R58 native source and its independently built companion accompany this manuscript; their identities are fixed in the source-provenance and preservation records.
\bibitem{Tzeng92} W.-G. Tzeng, \emph{A polynomial-time algorithm for the equivalence of probabilistic automata}, SIAM J. Comput. \textbf{21} (1992), 216--227. DOI: 10.1137/0221017.
\bibitem{LiYuYing12} Y. Li, N. Yu and M. Ying, \emph{Termination of nondeterministic quantum programs}, arXiv:1201.0891 (2012).
\bibitem{Fu24} J. Fu, H. Jiang, M. Xu, Y. Deng and Z.-B. Li, \emph{Algorithmic analysis of termination problems for nondeterministic quantum programs}, arXiv:2402.15827 (2024).

\bibitem{LiuVerdu18} J. Liu and S. Verd\'u, \emph{Rejection sampling and noncausal sampling under moment constraints}, Proc. IEEE International Symposium on Information Theory (2018), 1565--1569. DOI: 10.1109/ISIT.2018.8437857.
\bibitem{GocFlamich24} D. Goc and G. Flamich, \emph{On channel simulation with causal rejection samplers}, Proc. IEEE International Symposium on Information Theory (2024). DOI: 10.1109/ISIT57864.2024.10619339; examined text arXiv:2401.16579v2 (3 May 2024).
\bibitem{LiLiu26} Z. Li and T. Liu, \emph{Ranking experiments under sequential sampling}, arXiv:2608.19897v2 (28 August 2026). Version 1 appeared under the title \emph{The order of binary experiments under endogenous stopping}, with Z. Li as sole author; the version-2 statement is the one compared here.
\bibitem{Ying14pre} S. Ying and M. Ying, \emph{Reachability analysis of quantum Markov decision processes}, arXiv:1406.6146 (2014), Theorem 3.6; published journal version listed above as \cite{Ying18}.

\bibitem{KaraYuksel20} A. D. Kara and S. Y\"uksel, \emph{Robustness to incorrect system models in stochastic control}, SIAM J. Control Optim. \textbf{58} (2020), 1144--1182. DOI: 10.1137/18M1208058; examined text arXiv:1803.06046v7.
\bibitem{Demirci25} Y. E. Demirci, A. D. Kara and S. Y\"uksel, \emph{Sensitivity of filter kernels and robustness bounds to transition and measurement kernel perturbations in partially observable stochastic control}, arXiv:2508.10658v3 (2025).
\bibitem{Watrous18} J. Watrous, \emph{The Theory of Quantum Information}, Cambridge University Press, 2018, Chapter 3, Corollary 3.40.
\bibitem{KRN65} K. Kuratowski and C. Ryll-Nardzewski, \emph{A general theorem on selectors}, Bull. Acad. Polon. Sci. S\'er. Sci. Math. Astronom. Phys. \textbf{13} (1965), 397--403.

\bibitem{Charnes63} A. Charnes, W. W. Cooper and K. O. Kortanek, \emph{Duality in semi-infinite programs and some works of Haar and Carath\'eodory}, Management Sci. \textbf{9} (1963), 209--228. DOI: 10.1287/mnsc.9.2.209.
\bibitem{Shapiro09} A. Shapiro, \emph{Semi-infinite programming, duality, discretization and optimality conditions}, Optimization \textbf{58} (2009), 133--161. DOI: 10.1080/02331930902730070; author preprint, Optimization Online (2008).
\bibitem{Dufour25} F. Dufour and T. Prieto-Rumeau, \emph{Absorbing Markov decision processes and their occupation measures}, SIAM J. Control Optim. \textbf{63} (2025), 676--698. DOI: 10.1137/24M1661959. Related author preprint examined: \emph{A note on weak compactness of occupation measures for an absorbing Markov decision process}, arXiv:2402.10672v2 (2024).
\bibitem{Franc23} V. Franc, D. Prusa and V. Voracek, \emph{Optimal strategies for reject option classifiers}, J. Mach. Learn. Res. \textbf{24} (2023), no. 11, 1--49.
\bibitem{Jencova16} A. Jen\v{c}ov\'a, \emph{Comparison of quantum channels and statistical experiments}, arXiv:1512.07016v2 (2016). The separately indexed shorter version arXiv:1601.06370 was withdrawn in favor of this extended text.
\bibitem{LiElGamal18} C. T. Li and A. El Gamal, \emph{Strong functional representation lemma and applications to coding theorems}, IEEE Trans. Inform. Theory \textbf{64} (2018), 6967--6978; arXiv:1701.02827v4.
\bibitem{AldousShepp87} D. Aldous and L. Shepp, \emph{The least variable phase type distribution is Erlang}, Comm. Statist. Stochastic Models \textbf{3} (1987), 467--473. DOI: 10.1080/15326348708807067.
\bibitem{Sion58} M. Sion, \emph{On general minimax theorems}, Pacific J. Math. \textbf{8} (1958), 171--176. DOI: 10.2140/pjm.1958.8.171.
\bibitem{PiunovskiyZhang24} A. Piunovskiy and Y. Zhang, \emph{Extreme occupation measures in Markov decision processes with an absorbing state}, SIAM J. Control Optim. \textbf{62} (2024), 65--90. DOI: 10.1137/23M1572398. Author text examined: \emph{Extreme occupation measures in Markov decision processes with a cemetery}, arXiv:2307.03158v1.
\bibitem{DufourMixtures} F. Dufour and T. Prieto-Rumeau, \emph{Absorbing Markov decision processes: geometric properties and sufficiency of finite mixtures of deterministic policies}, arXiv:2512.17511; examined HTML at version-1 URL, Theorems 4.3--4.4 and Proposition 4.5.
\bibitem{Mereghetti24} C. Mereghetti, B. Palano and P. Raucci, \emph{Unary quantum finite state automata with control language}, Appl. Sci. \textbf{14} (2024), no. 4, 1490. DOI: 10.3390/app14041490.
\bibitem{RabinScott59} M. O. Rabin and D. Scott, \emph{Finite automata and their decision problems}, IBM J. Res. Develop. \textbf{3} (1959), 114--125. DOI: 10.1147/rd.32.0114.
\bibitem{Jaeger00} H. Jaeger, \emph{Observable operator models for discrete stochastic time series}, Neural Comput. \textbf{12} (2000), 1371--1398. DOI: 10.1162/089976600300015411.
\bibitem{ThonJaeger15} M. Thon and H. Jaeger, \emph{Links between multiplicity automata, observable operator models and predictive state representations---a unified learning framework}, J. Mach. Learn. Res. \textbf{16} (2015), 103--147.
\bibitem{MonrasWinter14} A. Monras and A. Winter, \emph{Quantum learning of classical stochastic processes: the completely-positive realization problem}, LIPIcs \textbf{27}, TQC 2014 (2014), 99--109. DOI: 10.4230/LIPIcs.TQC.2014.99.

\bibitem{Heller65} A. Heller, \emph{On stochastic processes derived from Markov chains}, Ann. Math. Statist. \textbf{36} (1965), 1286--1291. DOI: 10.1214/aoms/1177700000.
\bibitem{CarlylePaz71} J. W. Carlyle and A. Paz, \emph{Realizations by stochastic finite automata}, J. Comput. System Sci. \textbf{5} (1971), 26--40. DOI: 10.1016/S0022-0000(71)80005-3.
\bibitem{Fanizza24} M. Fanizza, J. Lumbreras and A. Winter, \emph{Quantum theory in finite dimension cannot explain every general process with finite memory}, Comm. Math. Phys. \textbf{405} (2024), article 50. DOI: 10.1007/s00220-023-04913-4; arXiv:2209.11225. See Section 2.3 for the FRDN phase examples and Theorem 2 and Section 4.2 for the noise-dependent worst-phase classical-memory lower bound.
\bibitem{DmitrievDynkin46} N. Dmitriev and E. Dynkin, \emph{On characteristic roots of stochastic matrices}, Izv. Akad. Nauk SSSR Ser. Mat. \textbf{10} (1946), no. 2, 167--184. Math-Net.Ru: im3595 (Russian).
\bibitem{JohnsonPaparella17} C. R. Johnson and P. Paparella, \emph{A matricial view of the Karpelevi\v c theorem}, Linear Algebra Appl. \textbf{520} (2017), 1--15. DOI: 10.1016/j.laa.2017.01.009; arXiv:1611.06970v2.
\bibitem{KirklandSmigoc22} S. Kirkland and H. \v Smigoc, \emph{Stochastic matrices realising the boundary of the Karpelevi\v c region}, Linear Algebra Appl. \textbf{635} (2022), 116--138. DOI: 10.1016/j.laa.2021.11.016; arXiv:2110.01040v1.
\bibitem{BenvenutiFarina04} L. Benvenuti and L. Farina, \emph{A tutorial on the positive realization problem}, IEEE Trans. Automat. Control \textbf{49} (2004), 651--664. DOI: 10.1109/TAC.2004.826715.
\bibitem{BenvenutiFarina99} L. Benvenuti and L. Farina, \emph{An example of how positivity may force realizations of ``large'' dimension}, Systems Control Lett. \textbf{36} (1999), 261--266. DOI: 10.1016/S0167-6911(98)00098-X.
\bibitem{OCinneide90} C. A. O'Cinneide, \emph{Characterization of phase-type distributions}, Comm. Statist. Stochastic Models \textbf{6} (1990), 1--57. DOI: 10.1080/15326349908807134.
\bibitem{Maier91} R. S. Maier, \emph{The algebraic construction of phase-type distributions}, Comm. Statist. Stochastic Models \textbf{7} (1991), 573--602. DOI: 10.1080/15326349108807207.
\bibitem{HoneVarona19} A. N. W. Hone and J. L. Varona, \emph{Continued fractions and irrationality exponents for modified Engel and Pierce series}, Monatsh. Math. \textbf{190} (2019), 501--516. DOI: 10.1007/s00605-018-1244-1. Author version arXiv:1810.10127v1, equation (3.3).
\bibitem{Nussbaum98} R. D. Nussbaum and M. Scheutzow, \emph{Admissible arrays and a nonlinear generalization of Perron--Frobenius theory}, J. London Math. Soc. (2) \textbf{58} (1998), 526--544. DOI: 10.1112/S0024610798006796.
\bibitem{Brown64} D.~R. Brown, On clans of non-negative matrices,
\emph{Proc. Amer. Math. Soc.} \textbf{15} (1964), 671--674,
doi:10.1090/S0002-9939-1964-0165504-6.
\bibitem{Schwarz64} \v{S}. Schwarz, On the structure of the semigroup of
stochastic matrices, \emph{Publ. Math. Inst. Hungar. Acad. Sci.}
\textbf{9} (1964), 297--311, especially Theorems~7--8.
\bibitem{Flor69} P. Flor, On groups of non-negative matrices,
\emph{Compositio Math.} \textbf{21} (1969), 376--382,
\url{https://www.numdam.org/article/CM_1969__21_4_376_0.pdf}.
\bibitem{Serre07} J.-P. Serre, Bounds for the orders of the finite subgroups
of $G(k)$, in \emph{Group Representation Theory}, M. Geck, D. Testerman
and J. Th\'evenaz (eds.), EPFL Press, Lausanne, 2007, 405--450;
author lecture text, April 2006, archived as arXiv:1011.0346.
\bibitem{BabaiBealsRockmore93} L. Babai, R. Beals and D. Rockmore,
Deciding finiteness of matrix groups in deterministic polynomial time,
\emph{Proceedings of ISSAC 1993}, 117--126,
doi:10.1145/164081.164104.
\bibitem{KieferWachter14} S. Kiefer and B. Wachter, Stability and complexity
of minimising probabilistic automata, in \emph{ICALP 2014, Part II},
Lecture Notes in Computer Science \textbf{8573}, 268--279,
doi:10.1007/978-3-662-43951-7\_23; arXiv:1404.6673, Section~4.2.
\bibitem{Schwarz62} \v{S}. Schwarz, Homomorphisms of a completely simple
semigroup onto a group, \emph{Matematicko-fyzik\'alny \v{c}asopis}
\textbf{12} (1962), no.~4, 293--300;
\url{https://www.dml.cz/handle/10338.dmlcz/126333}.

\bibitem{GPT13} J. Gouveia, P. A. Parrilo and R. R. Thomas,
\emph{Lifts of convex sets and cone factorizations}, Math. Oper. Res.
\textbf{38} (2013), 248--264. DOI: 10.1287/moor.1120.0575;
arXiv:1111.3164, especially Theorem 2.12 on equivariant factorizations.
\bibitem{KPT10} V. Kaibel, K. Pashkovich and D. O. Theis,
\emph{Symmetry matters for sizes of extended formulations},
SIAM J. Discrete Math. \textbf{26} (2012), 1361--1382.
DOI: 10.1137/110839813; arXiv:0911.3712v3.
An extended abstract appeared in IPCO 2010, LNCS 6080, 135--148.
\bibitem{BocciCarliniRapallo11} C. Bocci, E. Carlini and F. Rapallo,
\emph{Perturbation of matrices and nonnegative rank with a view toward
statistical models}, SIAM J. Matrix Anal. Appl. \textbf{32} (2011),
1500--1512. DOI: 10.1137/110825455; arXiv:1010.3566.
\bibitem{GillisGlineur12} N. Gillis and F. Glineur,
\emph{On the geometric interpretation of the nonnegative rank},
Linear Algebra Appl. \textbf{437} (2012), 2685--2712.
DOI: 10.1016/j.laa.2012.06.038; arXiv:1009.0880.
\bibitem{Okay26} C. Okay, \emph{Polyhedral classical simulators for quantum
computation}, Res. Math. Sci. \textbf{13} (2026), article 55.
DOI: 10.1007/s40687-026-00626-x. In particular Theorem 16, Lemma 17
and Corollary 19 treat invariant-polytope updates and adaptive composition.
\bibitem{AubrunSzarek17} G. Aubrun and S. J. Szarek,
\emph{Dvoretzky's theorem and the complexity of entanglement detection},
Discrete Anal. (2017), paper no. 1, 20 pp.; arXiv:1510.00578,
especially Section 6 on polyhedral approximations of density matrices.
\bibitem{Fawzi21} H. Fawzi,
\emph{On polyhedral approximations of the positive semidefinite cone},
Math. Oper. Res. \textbf{46} (2021), 1479--1489.
DOI: 10.1287/moor.2020.1077; arXiv:1811.09649.
\bibitem{HolroydLiggett16} A. E. Holroyd and T. M. Liggett,
\emph{Finitely dependent coloring}, Forum Math. Pi \textbf{4} (2016),
e9, 43 pp. DOI: 10.1017/fmp.2016.7; arXiv:1403.2448.
See Section 5 and the finite hidden-Markov obstruction attributed there
to Schramm.
\bibitem{Serre77} J.-P. Serre, \emph{Linear Representations of Finite Groups},
Graduate Texts in Mathematics \textbf{42}, Springer-Verlag, New York, 1977.
DOI: 10.1007/978-1-4684-9458-7.

\bibitem{MonrasWinter16} A. Monras and A. Winter,
\emph{Quantum learning of classical stochastic processes: the
completely-positive realization problem}, J. Math. Phys. \textbf{57}
(2016), 015219. DOI: 10.1063/1.4936935; arXiv:1412.3634.

\end{thebibliography}

\end{document}
